%  LaTeX support: latex@mdpi.com 
%  For support, please attach all files needed for compiling as well as the log file, and specify your operating system, LaTeX version, and LaTeX editor.

%=================================================================
\documentclass[coasts,article,accept,pdftex,oneauthor]{Definitions/mdpi} 

%=================================================================

% MDPI internal commands
\firstpage{1} 
\makeatletter 
\setcounter{page}{\@firstpage} 
\makeatother
\pubvolume{1}
\issuenum{1}
\articlenumber{0}
\pubyear{2024}
\copyrightyear{2024}
\externaleditor{Academic Editor: Firstname \linebreak Lastname}
\datereceived{25 December 2023} 
\daterevised{15 February 2024}
\dateaccepted{20 February 2024} 
\datepublished{} 
\hreflink{https://doi.org/} % If needed use \linebreak
%\doinum{}

%=================================================================
% Add packages and commands here. The following packages are loaded in our class file: fontenc, inputenc, calc, indentfirst, fancyhdr, graphicx, epstopdf, lastpage, ifthen, lineno, float, amsmath, setspace, enumitem, mathpazo, booktabs, titlesec, etoolbox, tabto, xcolor, soul, multirow, microtype, tikz, totcount, changepage, attrib, upgreek, cleveref, amsthm, hyphenat, natbib, hyperref, footmisc, url, geometry, newfloat, caption

\graphicspath{{figures/}} % path to folder with figures
%\usepackage{subcaption}
\usepackage[labelformat=simple]{subcaption}
\renewcommand\thesubfigure{\alph{subfigure}}
\DeclareCaptionLabelFormat{subcaptionlabel}{\normalfont(\textbf{#2}\normalfont)}
\captionsetup[subfigure]{labelformat=subcaptionlabel}



\usepackage{gensymb} % for \textdegree
\usepackage{tabularx}

%\usepackage{listings}
\usepackage{xcolor}
\usepackage[para,online,flushleft]{threeparttable}
\usepackage{array}
\usepackage{makecell, multirow}
\usepackage{booktabs}
\usepackage{caption}
\usepackage{adjustbox}
\usepackage{verbatim}





\usepackage{listings}
\lstset{xleftmargin=9.5pt,xrightmargin=3.5pt}
\captionsetup[lstlisting]{position=top, margin=0.75cm, labelfont={bf, small, stretch=1.17}, labelsep=period, textfont={small, stretch=1.17}, aboveskip=6pt, justification=justified}
%{\lstset{framextopmargin=2pt,framexbottommargin=2pt,abovecaptionskip=-3pt,
%belowcaptionskip=3pt,xleftmargin=1.3em,xrightmargin=4em}

%\captionsetup[lstlisting]{position=top, margin=0cm,labelfont={bf, small, stretch=1.17}, labelsep=period, textfont={small, stretch=1.17}, aboveskip=6pt, singlelinecheck=off, justification=justified}



\definecolor{deepblue}{rgb}{0,0,0.5}
\definecolor{deepred}{rgb}{0.6,0,0}
\definecolor{deepmagenta}{RGB}{195,12,214}
\definecolor{deepgreen}{RGB}{29,122,30}
\definecolor{mintgreen}{RGB}{192,249,192}
\definecolor{codepurple}{RGB}{204, 0, 153}
\definecolor{LightCyan}{rgb}{0.88,1,1}
\definecolor{GainsboroPurple}{RGB}{247,176,247}
\definecolor{LightSlateGrey}{RGB}{124,118,118}
%    
\lstdefinestyle{mystyle}{
    commentstyle=\color{LightSlateGrey},
    keywordstyle=\color{deepgreen}\bfseries,
    emphstyle=\color{deepred},
    numberstyle=\tiny\color{deepmagenta},
    stringstyle=\color{codepurple},
    basicstyle=\ttfamily\scriptsize,
    breakatwhitespace=false,         
    breaklines=true,                 
    captionpos=b,                    
    keepspaces=true,                 
    numbers=left,                    
    numbersep=5pt,                  
    showspaces=false,                
    showstringspaces=false,
    showtabs=false,                  
    tabsize=2
}
\lstset{style=mystyle}

%=================================================================
% Full title of the paper (Capitalized)
\Title{Random %MDPI: Title is different from the one submitted online at susy.mdpi.com. Please confirm which one is correct. %<-- Author's reply: the currently used is correct: 
 Forest Classifier Algorithm of Geographic Resources Analysis Support System Geographic Information System for Satellite Image Processing: Case Study of Bight of Sofala, Mozambique} %Attention: title altered. %<-- Author's reply: Ok.

% MDPI internal command: Title for citation in the left column
\TitleCitation{Random Forest Classifier Algorithm of Geographic Resources Analysis Support System Geographic Information System for Satellite Image Processing{: Case Study of Bight of Sofala, Mozambique}}

% Author Orchid ID: enter ID or remove command
\newcommand{\orcidauthorA}{0000-0002-5759-1089} % Polina Add \orcidA{} behind the author's name

% Authors, for the paper (add full first names)
\Author{Polina Lemenkova %MDPI: Please carefully check the accuracy of name and affiliation. %<-- Author's reply: yes, correct.
 \orcidA{}}

%\longauthorlist{yes}

% MDPI internal command: Authors, for metadata in PDF
\AuthorNames{Polina Lemenkova}

% MDPI internal command: Authors, for citation in the left column
\AuthorCitation{Lemenkova, P.}
% If this is a Chicago style journal: Lastname, Firstname, Firstname Lastname, and Firstname Lastname.

% Affiliations / Addresses (Add [1] after \address if there is only one affiliation.)
\address[1]{%
Department of Geoinformatics, Faculty of Digital and Analytical Sciences, Universität Salzburg, Schillerstraße 30, A-5020 Salzburg, Austria;  polina.lemenkova@plus.ac.at; Tel.: +43-677-6173-2772
}

% Contact information of the corresponding author
%\corres{Correspondence: polina.lemenkova@plus.ac.at; Tel.: +43-677-6173-2772 (P.L.)}

% Current address and/or shared authorship
%\firstnote{Current address: Affiliation 3.} 

% Abstract (Do not insert blank lines, i.e. \\) 
\abstract{Mapping coastal regions is important for environmental assessment and for monitoring spatio-temporal changes. Although {traditional} cartographic methods using \textcolor{black}{a geographic information system} (GIS) {are} applicable in image classification, \textcolor{black}{machine learning} (ML) methods present more advantageous solutions \textcolor{black}{for} pattern-finding tasks such as the automated detection of landscape patches in heterogeneous landscapes. This study \textcolor{black}{aimed} to discriminate {landscape} patterns along the eastern coasts of Mozambique using the ML modules of a Geographic Resources Analysis Support System (GRASS) GIS. The \textcolor{black}{random forest} (RF) algorithm of \textcolor{black}{the} module `r.learn.train' \textcolor{black}{was} used to map the coastal landscapes of \textcolor{black}{the} eastern shoreline of the Bight of Sofala, using remote sensing (RS) data at multiple temporal scales. The dataset \textcolor{black}{included} Landsat 8-9 OLI/TIRS imagery collected \textcolor{black}{in the} dry period during 2015, 2018, and 2023, {which \textcolor{black}{enabled the evaluation of} temporal dynamics}. The supervised {classification} of RS rasters was supported by the Scikit-Learn ML package of Python embedded in the GRASS GIS. The Bight of Sofala is \textcolor{black}{characterized} by \textcolor{black}{diverse} marine ecosystems dominated \textcolor{black}{by} swamp wetlands and mangrove forests located in the mixed \textcolor{black}{saline--fresh} waters along the eastern coast of Mozambique. %Please verify changes %<-- Author's reply:  yes, corrections are coloured blue. %<-- Author's reply: all corrections in abstract are now de-coloured to clean the text. 
 \textcolor{black}{This} paper demonstrates the advantages of using \textcolor{black}{ML} {for} RS data classification {in} the environmental monitoring of coastal areas. {The} integration of Earth Observation data, processed using \textcolor{black}{a} decision tree classifier \textcolor{black}{by} ML methods and land cover characteristics enabled \textcolor{black}{the detection of} recent changes in the coastal ecosystem of Mozambique, East Africa.
}

% Keywords
\keyword{machine learning; ensemble learning; GRASS GIS; Scikit-Learn; Python; random forest; satellite image; image processing; Africa; image classification} 

% The fields PACS, MSC, and JEL may be left empty or commented out if not applicable
\PACS{91.10.Da; 91.10.Jf; 91.10.Sp; 91.10.Xa; 96.25.Vt; 91.10.Fc; 95.40.+s; 95.75.Qr; 95.75.Rs; 42.68.Wt}
\MSC{94A08; 68U10; 54H30; 86A30; 86-08; 86A99}
\JEL{Y91; Q20; Q24; Q23; Q3; Q01; R11; O44; O13; Q5; Q51; Q55; N57; C6; C61}

\begin{document}

%----------- section ----------->
\section{Introduction}

%----------- subsection ----------->
\subsection{Background}
Remote sensing (RS) data {have} long been an important and efficient source of information {for} environmental monitoring, being one of the main data sources for mapping and detecting patterns {in} biodiversity across large spatial areas \cite{FAGAN2024432}. The~information from the Earth Observation satellite missions, which is provided quickly and {effectively}, %Please check your intended meaning is retained.  %<-- Author's reply:
has sufficient coverage{, }and {the} high quality of images {has} ensured numerous applications of RS data in modern geospatial studies and cartographic works~\cite{LOEB2023}. {Information derived from {satellite} images enables {the identification of} landscape characteristics that affect biodiversity {patterns}, the structural and functional properties of landscapes{, and the} spatial extent of different components of ecosystems. Such advantages of {RS} data enable {using} }both high-resolution and medium-resolution {products} in vegetation and land cover mapping{,} in order to detect patterns {in} vegetation changes{ }and~to track ecological interactions through time series analyses~\cite{6352393,info14040249,Lemenkova202324}. 

{Applications of RS data in environmental analyses include diverse} types of satellite images{,} {{which} can be used }for the effective interpretation of Earth's landscapes{,} {including} {Sentinel} 2 {data}~\cite{ZHAO2024169152,Lemenkova202021,CHOUDHARY20222443}, a National Oceanic and Atmospheric Administration Advanced Very High Resolution Radiometer~\cite{JI2017215,ZHAN2022112836}, Satellite Pour l'Observation de la Terre (SPOT{)}, \cite{FENSHOLT20091886}, {and}~Landsat~\cite{CHEN2023e19654,Lemenkova202323,LEMESIOS2024101153,SENAY2022113011,Lemenkova202321}, {as well as}~combinations thereof~\cite{XIE2023,RADELOFF2024113918}. For~instance, {the} recent Global Land Cover 2000 map was based on an analysis of the {vegetation} sensor on board the SPOT-4 satellite, using methods of digital image processing and GIS~\cite{GIRI2005123,WANG2023311}. {Moreover, RS data can be integrated with information about the environmental setting for analysis, mapping{,} and {modeling} purposes. For~instance, }Landsat scenes widely used for landscape mapping include the data{ }acquired from {its }various{ }sensors{,} such as {the} Multispectral Scanner (MSS), Thematic Mapper (TM), Enhanced Thematic Mapper Plus (ETM{+)}, and recent {products} of {the} Operational Land Imager and Thermal Infrared Sensor (OLI/TIRS), {or used {with} data integration }\cite{HE2018181,HUSSAIN2022103117,ESTOQUE2015205}.

The Landsat OLI-TIRS satellite missions {measured the} surface reflectivity and {scattering} in the {multi-spectral} bands{,} {which {resulted} in its successful application during {recent} decades} \cite{WULDER2022113195}. Its {capacity} to discriminate between different vegetation types {was} demonstrated in previous investigations~\cite{Ferreira,Montfort,Fatoyinbo,Ribeiro}. In~cloud-prone and rainy areas typical {of} coastal regions of tropical African countries, however, the~use of RS data is restricted by {atmospheric} effects{,} such as the attenuation of signals and cloud coverage during wet {periods}. {Therefore, an~analysis of vegetation should be based on cloud-free images or with {minimized} cloudiness (below 10}\%). {Moreover, the~remaining atmospheric effects should be {minimized} using advanced methods of computer vision. These include, for~instance, {atmospheric} scattering or haze{,} which occur as a result of{ }suspended dust {that accumulates} in relatively dry air. Such effects may lead to {color} distortion {of} the images and worsen image quality through impaired visibility{,} which may affect image analysis. }

As a result, the~processing of such data{ }requires advanced methods of satellite image processing{, e.g.,~removing noise, highlighting edges, {and} improving image contrast through the correction of atmospheric effects. This is possible using} {computer vision} and {machine learning} {(ML)} {algorithms} for data processing and visualization {due to {their} impressive computational and spatial analysis capabilities }\cite{LI2023104653,PHAM2023104501,HAN202387}. {The most {important} advantage of ML techniques is that {they are} capable of {automagically} deriving information from{ }existing datasets. Moreover, ML methods apply programming algorithms {that enable them} to excel in complex analyses and {modeling} of geospatial data} \cite{KHAN2022101678}. {In{ }environmental studies, valuable insights {from} ML {include} information extraction from complex terrain and {heterogeneous} landscapes{,} such as coastal regions. Such features of ML facilitate environmental analysis and {enhanced} planning and management of the coastal zones.}

Establishing new methods of Earth Observation data processing, such as {with} ML{,} requires the application of programming approaches {that} enable {the automation of} tasks~\cite{MULLISSA2023113799,Lemenkova202206}. Indeed, the~integrated use of{ }programming scripts, ML capabilities, and cartographic methods {has} provided powerful tools for RS data processing, with the aim of mapping, analyzing, and monitoring coastal areas using {the} distinctive sensor responses of land cover types. {Thus, {scripting} and programming {approaches enable overcoming} the ambiguities in image classification resulting from similarities {in the} spectral reflectance {of} various land cover types, especially with{ }medium-resolution Landsat data{, maintaining} the {high-automation} streams of {modeling}.} Advancements in computer processing; ML methods; and the development of scripting languages, such as Python, used either {alone} or {as }integrated {tools and add-ons }in a geographic information system (GIS), have allowed for an advanced {approach to} cartography and image processing. Such methods support accurate satellite data processing and analysis, which have {made} dynamic vegetation modeling possible~\cite{Lemenkova202227}. 

%----------- subsection ----------->
\subsection{Related~Works}
Generally speaking, the~accuracy of classification of {the} vegetation types in heterogeneous landscapes using satellite images {mainly depends} on {the} sensitivity of sensor backscattering to{ }{differences} between the phenological characteristics of plant leaves{. Alternatively, they should also be sensitive to the variations in} structure of {the} other land cover types{ }such as bare land or urban areas with dominating impervious surfaces{. Hence, }different {interactions} between the sensor backscatter and the structure of the canopy of landscape patches {can provide {novel} information on ecological structure}. Nevertheless, many recent studies~\cite{rs15235559,WU2023103407,CHEN2023101499,rs15225387,LIN2021102370}{ }noted the evident advantages of using {ML} methods in image processing. For~instance, ref.~\cite{LIAO2023110231} found that the use of deep learning (DL) modeling in remote sensing data processing could improve {the accuracy of the }classification. 

Earlier studies {have} also reported that {the} classification precision {was significantly} improved {when} the{ }computer vision algorithms of DL for image analysis and interpretation {were applied}~\cite{GUO20231,rs15215121,DOU2021102477}. Similar studies also validated the potential of the use of ML methods in cartographic{ }applications {of RS data processing }for mapping various vegetation classes{. Thus, these methods were used for evaluation of }the optimal combination of geospatial data {used }for separating different land cover classes~\cite{BOULILA2021106014,LIN2024102618}. These and other{ }results {reported in previous studies }also indicated that{ }the use of Python {improves the performance of} the classification and {attains} the required automation and accuracy level in {the process of }image classification~\cite{ROBERTS2022105192,GUO2023e21542,Redoloza}. Nevertheless, improvements and flexibility {are} achieved in the Geographic Resources Analysis Support System (GRASS) GIS software{, which} presents a smart combination of {the existing }Python libraries with cartographic toolsets and {a} powerful image processing~functionality. 

{Since the {release of the }first version of{ }GRASS GIS{ }in 1984} \cite{GRASS_GIS_software}, {{this} instrument has been constantly improved, which resulted in a {recent }updated version of software 8.3.1{. Nowadays, GRASS GIS presents} a powerful geospatial data processing engine {with an} integrated suite of {the }existing modules for raster and vector data processing. {The }advances in technical progress and the rapid development of cartographic instruments along with the {improvements} in programming algorithms resulted in various types of modules of GRASS GIS that include diverse geospatial data processing tools. The~applications of such tools {can be found} in a variety of case studies. To~mention a few of them, the~use of GRASS GIS in geographic studies includes environmental monitoring}~\cite{STRIGARO201668,Lemenkova202322}, {time series and massive data analysis} \cite{JASIEWICZ20111525,Lemenkova202313}, {hydrological modeling} \cite{JASIEWICZ20111162}, {computing landscape diversity, image segmentation} \cite{Lemenkova202320,ROCCHINI201382}, {and geomorphometric modeling} \cite{HOFIERKA2009387}. {These{ }examples{ }prove {the} effectiveness {of the GRASS GIS }for image processing and remote sensing data analysis.}

%----------- subsection ----------->
\subsection{Objectives and~Motivation}

{The primary objective of the present research is to determine the spatio-temporal changes in the land cover types of coastal Mozambique through the application of machine learning methods to remote sensing data analyses. }In view of the advantages of the ML methods briefly discussed above, {this study applies} such methods {and presents} the case of image processing using the GRASS GIS module  r.learn.train for image classification. Specifically,{ }the random forest (RF) algorithm of the ML method{ is employed in this study }with the sample coastal area of eastern Mozambique, Bight of Sofala. Using ML algorithms embedded in the modules of GRASS GIS enables us to perform a supervised classification using training pixels of raster tiles and the integrated Python Scikit-Learn library~\cite{Pedregosa} for modeling vegetation in coastal~areas. 

{The goal of this study is to develop a workflow using several modules of GRASS GIS for ML-based random forest techniques of image classification. To~this end, the~multi-spectral Landsat {8--9} %MDPI: Please check the full text and confirm if this should be en dah ("--" U+2013). %<-- Author's reply: yes, corrected to dash
 OLI/TIRS images are used for land cover mapping of the coastal region of eastern Mozambique that encompasses the Bight of Sofala and cover a time gap from 2015 to 2023. Since the area is located in a cloud-prone and rainy region affected by occasional floods, the~data are selected during the `dry' period from August to September using optimal atmospheric conditions. Additionally, the~images are corrected to improve {their} quality. In~total, three scenes are used for classification and ML-based processing.}

{The main aim of this study is to evaluate short-term temporal dynamics in the region of the Bight of Sofala using advanced methods of cartographic scripting using the ML modules of {the} GRASS GIS. The~methodology includes the extraction of training data from{ }raster images, supervised ML methods that use ensemble classification{, }the regression tree method using RF, and cross-validation tasks. Technically, this paper demonstrates the advantages of using ML algorithms for RS data classification{, as well as advanced approaches of} cartographic mapping for {the }environmental monitoring of coastal areas. High robustness, versatility, and applicability of the RF algorithm motivated its application for processing satellite images in the shoreline regions {which are notable for highly vulnerable} environmental parameters{. }Due to the sensitivity of RF to noise and outliers in the pixel matrix, its application to image processing is beneficial for reducing {the} overfitting by averaging multiple decision trees{, as demonstrated} in this study. {As a result, the use of RF algorithm enabled us to reduce} the data processing time and {to increase the} accuracy of {the }Landsat {8--9} OLI/TIRS imagery. }

%----------- subsection ----------->
%\pagebreak
\section{Study~Area}

This study {is located in }the coastal {area} of Mozambique{, }western {part} of the Indian Ocean{{. Specifically, it investigates the region of the} Bight of Sofala, where the estuary of the Beira river enters the Mozambique Channel} (Figure~\ref{fig01}). 
 
\begin{figure}[H]
%	\centering
%		\hspace{-35pt}
		\includegraphics[width=10.0 cm]{Fig_01.jpg}
\caption{Study %MDPI: Please change the hyphen (-) into a minus sign ($-$, "U+2212"). e.g., "-1" should be "$-$1". % Author's reply: the map in figure 1 is corrected as requested (the hyphen (-) is changed into a minus sign ($-$, "U+2212")).
 area in the coastal region of the Bight of Sofala, Mozambique is indicated by the cyan-colored rotated square. Mapping software: Generic Mapping Tools (GMT). Map source:~author.}\label{fig01}
\end{figure}
 
{As} one of the developing countries of Africa, Mozambique is notable for its rapid land cover changes{ }related to the environmental effects and urbanization{. As a result, }landscape dynamics have {recently }become one of the major issues of the country{, }which is reported in many {relevant} studies~\cite{BACAR2023,JANSEN2008308,MUCOVA2018e00447}. Located on the southeast coast of Africa, Mozambique represents a country with {diverse topographic forms} and landscape setting{. The most remarkable geomorphic features} include a complex mix of {the} inland hills and low plateaus {situated }along the narrow coastal strip {which contrast with} rugged highlands {located }in the western parts of the country. {Such} variability in landscapes {favors diversified }vegetation {structures, increases {the }biodiversity} richness, and {creates }{unique} land cover {patterns} {that are }characteristic of{ Mozambique} (Figure~\ref{fig02}). 

 \begin{figure}[H]
	
\begin{adjustwidth}{-\extralength}{0cm}
%\centering %% If there is a figure in wide page, please release command \centering
	\centering
		\hspace{70pt}\includegraphics[width=15.5 cm]{Fig_02.jpg}
		\end{adjustwidth}
\caption{Land %MDPI: Please change the hyphen (-) into a minus sign ($-$, "U+2212"). e.g., "-1" should be "$-$1". % Author's reply: corrected: the hyphen (-) is changed into a minus sign ($-$, "U+2212")
 cover types in Mozambique according to the Food and Agriculture Organization (FAO) data. Mapping software: QGIS. Map source: author's~work.}\label{fig02}
\end{figure}

The major geomorphic types {of Mozambique }include numerous highlands interspersed with plateaus; covered with woodlands; and~lowlands, {which are }situated mostly to the south of the Zambezi River{---}the major river of {the country}. {The }Zambezi {River }roughly {separates} the country into the northern and southern parts, {which have notably }distinct ecosystems. Being the fourth-longest river in Africa with a complex hydrological setting, it largely affects the distribution of plants and species in the surrounding landscapes and {in the }estuary zone{, which is }located to the north of the Bight of Sofala~\cite{Dunham,Dunham1991,Simasiku}. The~hydrology of {the} Zambezi River is connected to the main channel and tributaries, and {water management involves} the eight countries surrounding the basin~\cite{Mbumwae}. %Please check your intended meaning is retained. 
{Thus}, through the downstream and upstream {movement} %Please check meaning is retained. 
of water that changes the local salinity and suspended sediments, the~hydrology of Zambezi strongly influences the coastal ecosystems of the Sofala Bank~\cite{NEHAMA2021101891}. {Nevertheless, the recent construction of artificial} dams has had significant environmental effects on the major floodplains{. For instance, these include} {reduced} water supply and river discharge~\cite{Gope} {and increased} sedimentation of the coastal landscapes of its estuary~\cite{Moore}. 

{The coastal ecosystems in the} Bight of Sofala {are} one of the {most} environmentally vulnerable regions {in} Mozambique. This is caused by{ }multiple impacts {from} the climate {setting }and {tidal hydrology} \cite{LEAL2009605} {including tidal} cycles, shelf circulation, and coastal currents{ }~\cite{MALAUENE20181}. In~turn, the~hydrological setting of the coastal ecosystems influence {the patterns of }primary production and the  distribution of zooplankton, phytoplankton, and nutrients in the {shelf} waters~\cite{MALAUENE2021107268}. Moreover, the rapid expansion of {urban} areas {caused by} population {growth }and economic pressure{ changes the landscapes}{. Such social processes }{result} in {urban sprawl, which changes the }natural landscapes {into artificial spaces }\cite{JOKARARSANJANI201838}. Moreover, the~Sofala Bank is known for its industrial shallow-water shrimp fishery~\cite{SOUSA2006207,MIGUEL2003365,GAMMELSROD199291}, which creates additional pressure on the coastal ecosystems through{ }the overexploitation of natural resources. 

{The}~repeated and severe floods and cyclones in{ }western part of the Indian Ocean have a great influence on{ }the landscapes and {distribution of }vegetation {along the coasts of} Mozambique. As~a result, Mozambique is reported to be one of the most vulnerable south African countries{{ }to} flooding and tropical cyclones{. Such climatic hazards }cause{ }damage to {the }infrastructure and {create additional factors to the disruption of }ecosystems~\cite{NHANGUMBE2023101015,SALVUCCI2020106713,Gall}. {Another{ }}environmental {issue in} the coastal {regions} of Mozambique {concerns }the decline of mangroves, {{which} is caused by the cumulative effects from }anthropogenic activities and climate change~\cite{COME2023106813}. {At the same time, }{as }highly productive coastal ecosystems with unique fauna and flora{, }mangrove swamps are important sources of natural resources and livelihoods for local {people}. Therefore, their decline and degradation negatively affect the sustainability of coastal ecosystems{.}

{Over{ }recent decades, the~}consequences of floods in Mozambique include the destruction of landscape structures~\cite{FICHTNER2023103329}{ }{and} crops during the flood{ }{period}~\cite{BOFANA2022112808}. {Moreover, this increased }landscape fragmentation {resulted in the continuous physical disintegration of habitats into smaller patches with decreased}{ }size and isolation {between{ }habitat clusters. The~consequences of such processes negatively affect the }biodiversity and species richness~\cite{Guldemond}{ }as~well as lead to disruptions in the coastal ecosystems~\cite{RAMAYANTI20221025}. In~terms of social--economic aspects of Mozambique, climate-related natural hazards and disrupted landscapes affect the farming sector, fishery{, }and~agriculture activities in the coastal {areas} of~Mozambique. 

All these factors {are required for an effective} environmental monitoring system{. Specifically for }the coastal regions of Mozambique{, a data analysis enables us }to define the cost of landscape restoration in the affected areas and {to }{detect} land cover{ }changes.{ }{Proper land} planning in {the }coastal areas of Mozambique requires {the costly process of }evaluating {landscape }dynamics{, }{performing a climate change} {assessment}, and carrying out an {integrated} socio-environmental {analysis}~\cite{SMITH2019101976,LUNDGREN2022103023,SILVA201843,MARTINS2022102793}.{ The }lowest possible cost when {monitoring} the coastal areas of Mozambique {can} be achieved by evaluating the difference between{ }{land cover} types and performing a patch analysis using the image classification and quantification of landscape fragments~\cite{PITTMAN2023}. This is possible using cartographic mapping and a time series analysis of the satellite images~\cite{XI2019111212,ABIDI2023106152,MOHAMMADI2023272}. Land surface topography, which varies significantly in Mozambique, influences the distribution of vegetation types that can be considered in agricultural planning, environmental management, {and} ecological monitoring of the coastal~regions. For~instance, to~preserve{ }the vulnerable ecosystem of the narrow strip along the coasts of the Bight of Sofala, operative monitoring using remote sensing data {is }a useful technical tool~\cite{NAMAGANDA2023117811}. 

%----------- section ----------->
\section{Materials and~Methods}
\unskip

\subsection{Workflow}
In this paper, multi-temporal Landsat {8--9} OLI/TIRS images {were} used for land cover mapping in the coastal region of the Bight of Sofala, eastern Mozambique,{ }Indian Ocean. {The{ }methodology included the {extraction of the} training data from{ }raster images, supervised ML using ensemble classification{, }the regression tree method using RF, and cross-validation tasks}. {A schematic view of the workflow in cartographic and remote sensing {tasks} is portrayed in Figure} \ref{fig03}{.}

\vspace{-6pt} \begin{figure}[H]
	
\begin{adjustwidth}{-\extralength}{0cm}
%\centering %% If there is a figure in wide page, please release command \centering
	\centering
		\hspace{70pt}\includegraphics[width=15.5 cm]{Fig_03.jpg}
\end{adjustwidth}
\caption{{Workflow of the data processing. Diagram software: R. Source: author's work.}}\label{fig03}
\end{figure}

The workflow was {completed using GRASS GIS software and} {includes} the following {14} steps:

\begin{enumerate}
  \item Image importation and interpretation {using the Geospatial Data Abstraction Library (GDAL){ }and reprojection of images{: }Since all the images must be registered in the software {and be coherent with each other},  %Please check intended meaning is retained. 
  the~registration procedure of the Landsat images was performed using the} `r.import' module. The~metadata of the satellite images {were  set to the current region and checked} %Please check intended meaning is retained. 
to assess the suitability and quality control: data acquisition, cloudiness{, }sun elevation {and {geographic} extent. Technically, this was carried out using {the }code ``r.import input{=}%MDPI: Please check the full text and confirm if space is needed around equal signs. % Author's reply: No, not needed.
/Users/<path>/ImageName.TIF output=ImageName extent=region resolution=region''. Here the path to the original image is provided in full; the resolution was set to the 30 m as for Landsat scenes}.
  \item Image preprocessing is a fundamental approach in remote sensing data analyses{. It  aims} to prepare the images for further processing and analyses. Here, the Landsat images were integrated with GDAL and preprocessed using the `i.landsat.toar'{ }module, which calculates the top-of-atmosphere spectral reflectance and temperature for Landsat OLI/TIRS scenes. To~this end, the~information on sun elevation and the date of image acquisition {were} used from the metadata (product date). The~digital numbers (DN) of the image pixels were converted to the top-of-atmosphere radiances (ToAR). {The resulting output images {were }processed using the `i.landsat.toar' module {and}  {the }DN of {the} images were transformed, which{ }calibrated {the} top-of-atmosphere~reflectance.}
  \item {To analyze the land cover types visible on the Earth's surface as a representation landscape patches, color} composites {were created using bands of the satellite images}. The~image composites were generated using `r.composite' for natural and false color composites using triplets of red, green, and blue bands merged into a single composite raster map {by the following code (the example is for the band combination {7-5-3}): ``r.composite blue=L807 green=L805 red=L803 output=L8753 {--overwrite}%MDPI: Please check the full text and confirm if space should be removed before en dash.
''. Here, each band was selected for each of the RGB triplets (red, green, and blue), which enabled us to visualize the images in various combinations (natural and false color composites)}. 
  \item {The landscapes of coastal Mozambique are characterized by a complex mosaic of heterogeneous land cover types that{ }exhibit structural similarities in spectral reflectance. Therefore, a }radiometric {partition of the images was performed} using the unsupervised k-means method {with the }`i.cluster' module {which employs computer vision techniques for objective and automatic recognition of land cover types. This method generates spectral signatures for different land cover types using a clustering algorithm that identifies the differences between various values of spectral reflectance for each pixels and assigns {them} to classes accordingly.}
    \item The generation of signature files for land cover types {was performed }using {the }'signaturefile' function of the `i.cluster' module. {The profiles of the }spectral backscatter coefficient{ }of the land cover types were generated and compared to {the }FAO classification with 10 major {identified} land cover classes{{ }in {the }coastal regions of Mozambique}. {Various} characteristics of landscape patches were identified and compared for{ }images taken at different {periods as a time series}{.}
  \item Image classification {{was performed to identify }land cover classes }using the maximum likelihood method (`i.maxlike' module). {This method was selected due to its versatility and ability to produce fast and accurate results. In~essence, the~MaxLike classification algorithm uses clusters generated in the previous step and classifies pixels{ }in an iterative self-organizing process of image partition to produce a set of spectral classes. {Afterwards,}  labels {are assigned }to the clusters automatically and{ }maps {are generated }using a defined color palette}.
  \item The generation of training pixels from an older land cover classification using the `r.random' module {was performed to generate }a set of raster point maps. {This module creates a raster map which contains coordinates of points whose locations have been randomly determined to ensure {the }objectivity and randomness of the representative set of pixels.{ }It } includes randomly located cells and pixel points using {a} non-deterministic random seed. {Technically, it was implemented by the following code: ``r.random input=L8CL seed=100 npoints=1000 raster=L8CLroi''.}
  \item The creation of an imagery group with all Landsat OLI/TIRS {multispectral }bands {that have 30 m resolution{, }excluding panchromatic bands: this was performed }using `i.group' module {with the following code: ``i.group group=L82015 input=L801,<...>,L807''}. 
  \item A random forest classification model was trained using {the }`r.learn.train' approach of {the }ML modules of {the }GRASS GIS. {Practically, it was implemented using the following code: r.learn.train group=L82015 training\_map=L82015roi \\
  model\_name=RandomForestClassifier n\_estimators=500 save\_model=rf\_model.gz}
  \item Performing predictions using the `r.learn.predict' module of GRASS GIS{: This is an essential part of the ML data processing }that uses a fitted Scikit-Learn estimator {derived} from Python to raster layers in a group of images. {The implementation was performed using the code ``r.learn.predict group=L8\_2015 load\_model=rf\_model.gz output=rf\_classification --overwrite''. {The} algorithm uses{ }the training {dataset} generated in {a }previous step and {applies} the created model ``rf\_model.gz''. {The}~main goal{ }is to apply a fitted estimator from the Python's ML library Scikit-Learn to {the imagery group }which was generated during the step ``i.group''.} 
  \item {Checking} raster categories using {the }module `r.category'{: }{The category values and labels {are displayed }for{ }the raster layer requested for the generated images{. This is carried out }by the command ``r.category rf\_classification'', {which} reports the data in the console output. In} this case, 10 major land cover classes of the coastal area of Mozambique were identified by the computer vision, which were automatically applied to the classification results.
  \item Visualization of the results {was performed }using {the GRASS GIS }module `r.colors', which copies {the }color scheme from {the }land class training map to the output maps{. For~instance, the~code for the maps of random forest classification {was} as follows: ``r.colors rf\_classification color=plasma-e''}
  \item The results {were mapped }using a set of modules---`d.rast', `d.legend', and `d.mon'---{that represent the cartographic tools of GRASS GIS {for mapping and the visual display of data}}. %Please check meaning
  The~files were {then }saved using the module `d.out.file' as bitmap graphics. {For instance, data visualization included the following snippets of code: ``d.rast rf\_classification d.legend raster=rf\_classification title=``Random Forest: 2018'' title\_fontsize=14 font=``Helvetica'' fontsize=12 bgcolor=white border\_color=white''}
  \item {Interpretation} of the results {was based on the} comparison of maps of land cover types of the Bight of Sofala, Eastern Mozambique, for 2015, 2018, and 2023 {to analyze the spatio-temporal dynamics of land cover types in the coastal landscapes}. 
\end{enumerate}  

%----------- subsection ----------->
\subsection{Data}

In this study, three multi-temporal images of the Landsat Operational Land Imager and Thermal Infrared Sensor (OLI/TIRS) were selected {to produce a classification {using }machine learning{. }The~images were chosen }due to the high quality of the data, availability {of the{ }high-quality 'leaf-on' data on coastal Mozambique, }and coverage of {the study area}. The~{most essential }technical characteristics of the images are summarized in Table~\ref{tab01}. 

\begin{table}[H]
\caption{Identifiers (ID) of the six Landsat 8--9 OLI/TIRS images on {the} study area {of Mozambique}, obtained from the {EarthExplorer} USGS~repository.}
\label{tab01}
\footnotesize
\begin{tabularx}{\textwidth}{p{2.6cm}ll}
			\toprule
			\textbf{Date} & \textbf{Landsat Product Identifier L1} & \textbf{Scene ID} \\
			\midrule
			9 July 2015 & \makecell{LC08\_L2SP\_167074\_20150709\_20200909\_02\_T1} & LC81670742015190LGN01 \\
			19 September 2018 & \makecell{LC08\_L2SP\_167074\_20180919\_20200830\_02\_T1} & LC81670742018262LGN00 \\
			24 August 2023 & \makecell{LC09\_L2SP\_167074\_20230824\_20230826\_02\_T1} & LC91670742023236LGN00 \\	
			\bottomrule
		\end{tabularx}
\end{table}

The images are presented in Figure~\ref{fig04}. Low cloudiness and acceptable resolution depend on the acquisition techniques of Landsat ensured by the United States Geological Survey (USGS). 

\begin{figure}[H]


\hspace{15pt}\makebox[0.90\linewidth][r]{%
\captionsetup[subfigure]{justification=centering}

	\begin{subfigure}[b]{.40\textwidth}
		\centering
			\includegraphics[width=0.9\textwidth]{Fig_04a.jpg}
		\caption{2015}
	\end{subfigure}%
	\begin{subfigure}[b]{.40\textwidth}
		\centering
		\includegraphics[width=0.9\textwidth]{Fig_04b.jpg}
		\caption{2018}
	\end{subfigure}%
	\begin{subfigure}[b]{.40\textwidth}
		\centering
		\includegraphics[width=0.9\textwidth]{Fig_04c.jpg}
		\caption{2023}
	\end{subfigure}%
}

\caption{Landsat images on the coastal region of Beira, Mozambique: (\textbf{a}) 9 July 2015; (\textbf{b}) 19~September 2018; (\textbf{c}) 24 August 2023. The~images show natural colors of the Landsat 8--9 OLI/TIRS.}\label{fig04}
\end{figure}

The {values of the }cloud {coverage were} 0.01, 0.11, and 0.69 for the images from 2015, 2018, and 2023, respectively. The~original satellite images are stored in Universal Transverse Mercator (UTM) Zone 36 for Mozambique. The~terrain data of the study area used for topographic visualization are based on the General Bathymetric Chart of the Oceans (GEBCO) grid with high spatial resolution and accuracy (15 arc seconds) and visualized as the digital terrain model Shuttle Radar Topography Mission (SRTM). The~topographic names are added {according} to the gazetteer of Digital Chart of the World (DCW).  

\subsection{Software}

Traditional image processing methods rely on the use of GIS~\cite{GEMUSSE2023101022,LIANG2023113367,BEY2020111611}. However, such methods have some drawbacks in that{ }the human-controlled maintenance of workflow is required for{ }image{ }processing, classification, and analysis, which might be hard{ }in {case of} {image series}. In~contrast, script-based image processing {uses} the automation of data processing {to derive} information {from the} RS data~\cite{Lemenkova202213}. An~alternative approach is presented by machine learning (ML),{ }which can be used as a  descriptor{ }of pixel values{ }and partial replacement of the traditional methods using programming~\cite{jimaging9050098,Lemenkova202229}. The~ML{ }is based on {the} algorithms of artificial intelligence, which include{ }computer vision and statistical {analyses}. {This} study applies ML{ using} the latest available version of GRASS GIS 8.3.1{.}

{The}~installation of the program was performed using {the} `sudo port install grass' with compilations, {which includes }the{ }dependencies {of} the essential Python's libraries and{ }ML packages. {The }programming is {implemented} in {the }GRASS GIS as scripts{. This automates} the repeated steps of image processing {to improve} the workflow~\cite{technologies11020046,jmse11040871}. The~data processing was performed on the MacBook Apple Sonoma (arm64 {architecture}). The~Cooperative Computing Tools (CCTools), which enable large-scale distributed computations from clusters, were updated and installed using {the }sudo {command} as follows: `sudo port clean cctools' and `sudo port {-v} %MDPI:  Please check the full text and confirm if space should be removed before hyphen. % Author's reply: No, space should be kept: it is a syntax of programming. 
 install cctools'. After~the CCTools and XCode were installed and activated, {all the }necessary updates in the libSystem of MacPorts were performed including the deactivated library libunwind-headers: `sudo port {-f} deactivate~libunwind-headers'. 

To employ the enhanced properties and improved functionality of GRASS GIS adjusted for satellite image processing, the~Python-based modules and the toolchain GNU Compiler Collection (GCC) {were used}. The~GCC 13 was installed accordingly for code compilation, supporting library dependencies through linking and conversion of GRASS GIS scripts into the binary executable format for assembly into  executable files. The~wxPython was installed additionally using the `sudo port install py311-wxpython-4.0' {command called} from the~MacPorts. The topographic map of the study area was made using the Generic Mapping Tools (GMT) software developed by~\cite{Wessel}, {which represents a toolbox for geospatial data processing}. The~method{ }used for mapping {includes} scripts {that }were derived from earlier works~\cite{Lemenkova202217,Lemenkova202212,Lemenkova202203}. 

\subsection{Scripting}
The advanced script-based cartographic ML methods of {the }GRASS GIS for monitoring and mapping land cover types were built due to the growing demand for time- and cost-effective approaches in the cartographic workflow and remote sensing data analyses{. It enables} users to save time and resources while processing {the }geospatial datasets and~{to }use open-source satellite images. {Three} major script-based toolsets were used in this study for the geospatial data analysis and environmental monitoring of West Africa:

 \begin{enumerate}
        \item GRASS GIS module `r.import' was used for importing the data, Listing \ref{lst01}; 
        \item GRASS GIS module `r.composite' was used for creating false and natural color composites, Listing \ref{lst02}; 
        \item GRASS GIS {modules} of `i.cluster' and `i.maxlik'  were used for unsupervised image classification and accuracy assessment {by} computed rejection probability classes for classified pixels, Listing \ref{lst03};  
         \item A combination of GRASS GIS modules `d.rast', `d.legend', `r.colors', `d.mon', `g.region', and `d.out.file' was used for mapping the output images, Listing \ref{lst04}; 
        \item GRASS GIS {module} `r.learn.train' was used for random forest classification of machine learning (ML)-based supervised image classification using the training data generated in a previous step, Listing \ref{lst05}. 
  \end{enumerate}
  
The main idea {of this framework }is to {perform} in a single workflow{ }both{ }image processing and{ }mapping using scripts run from the console. This enables a user to reveal the technical performance of {the }automated geospatial data processing using a programming approach with three different programming suits {of} the script-based software. Such an approach demonstrated{ }that the inclusion of scripts speeds up cartographic plotting and image processing and {facilitates} the research methodology. As~mentioned earlier~\cite{Lemenkova202215}, the~main advantage of scripts consists in automation of the process and repeatability of the workflow,{ which saves} time. {Using} the information received from the automated image analysis and interpretation{, }we were able to detect the changes in the analysis of the landscapes of eastern Mozambique for evaluation of the environmental effects from climate change{. This was achieved }through a {time }series {analysis }of the {satellite }images{ }taken for different years (2015--2023). 

First, the~GRASS GIS was started from the console, and then, the bands of the Landsat-8 OLIimage on 2015 were imported stepwise using the `r.import' module, as shown in the script of Listing \ref{lst01}. The~same algorithm was repeated for the Landsat images for 2018 and~2023.

\begin{listing}[H]
\caption{GRASS %MDPI: Please confirm if listings should be moved behind their first citations. % Author's reply: yes, confirm.
 GIS code for importing satellite images.}
\begin{lstlisting}[language=bash,caption={},style=mystyle,label={lst01}]
grass
r.import input=/Users/polinalemenkova/grassdata/Mozambique/LC08_L2SP_167074_20150709_20200909_02_T1_SR_B1.TIF output=L8_2015_01 extent=region resolution=region
r.import input=/Users/polinalemenkova/grassdata/Mozambique/LC08_L2SP_167074_20150709_20200909_02_T1_SR_B2.TIF output=L8_2015_02 extent=region resolution=region
r.import input=/Users/polinalemenkova/grassdata/Mozambique/LC08_L2SP_167074_20150709_20200909_02_T1_SR_B3.TIF output=L8_2015_03 extent=region resolution=region
r.import input=/Users/polinalemenkova/grassdata/Mozambique/LC08_L2SP_167074_20150709_20200909_02_T1_SR_B4.TIF output=L8_2015_04 extent=region resolution=region
r.import input=/Users/polinalemenkova/grassdata/Mozambique/LC08_L2SP_167074_20150709_20200909_02_T1_SR_B5.TIF output=L8_2015_05 extent=region resolution=region
r.import input=/Users/polinalemenkova/grassdata/Mozambique/LC08_L2SP_167074_20150709_20200909_02_T1_SR_B6.TIF output=L8_2015_06 extent=region resolution=region
r.import input=/Users/polinalemenkova/grassdata/Mozambique/LC08_L2SP_167074_20150709_20200909_02_T1_SR_B7.TIF output=L8_2015_07 extent=region resolution=region
g.list rast
\end{lstlisting}
\end{listing}\vspace{-12pt}

Second, {the }color composites were created using Listing \ref{lst02} for false and natural color composites with the case of an image from 2015. {While true color composites show a representation of the Earth visible naturally to human eyes, false color composites better display selected features on the Earth (e.g., discrimination between water and land surfaces). {This is done }through a combination of visible red, green, and blue bands that correspond to the multispectral channels of the images}.

\vspace{-6pt}
\begin{listing}[H]
\caption{GRASS GIS code for creating color composites.}
\begin{lstlisting}[language=bash,caption={},style=mystyle,label={lst02}]
r.composite blue=L8_2015_07 green=L8_2015_05 red=L8_2015_03 output=L8_2015_753 --overwrite
d.mon wx0
d.rast L8_2015_753
d.out.file output=Mozambique_753 format=jpg --overwrite
# true color
r.composite blue=L8_2015_02 green=L8_2015_03 red=L8_2015_04 output=L8_2015_234 --overwrite
d.mon wx0
d.rast L8_2015_234
d.out.file output=Mozambique_234 format=jpg --overwrite
# false color: NIR band B05 in the red channel, red band B04 in the green channel and green band B03 in the blue channel
r.composite blue=L8_2015_03 green=L8_2015_04 red=L8_2015_05 output=L8_2015_345 --overwrite
d.mon wx0
d.rast L8_2015_345
d.out.file output=Mozambique_345 format=jpg --overwrite
\end{lstlisting}
\end{listing}\vspace{-9pt}

Third, the~next step included clustering the image using the `i.cluster' module. The~computational region was set to the band of the Landsat image to match the scenes {of the satellite images}. The~bands of {the }Landsat channels were grouped using {the} `i.group' module, {which is designed }for grouping the data{. {The}~selected bands} were used as grouped data during {the }clustering {process. }The~signature file was generated using {the }'signaturefile' command. The~image was classified into {the }10 major classes of the classification scheme using {the} information from FAO. {The}~unsupervised automated classification was performed using the `i.maxlik' module. The~complete script {of this process }is shown in Listing \ref{lst03}. 
 
\begin{listing}[H]
\caption{GRASS GIS code for grouping bands and clustering.} 
\begin{lstlisting}[language=bash,caption={},style=mystyle,label={lst03}]
g.region raster=L8_2015_01 -p
i.group group=L8_2015 subgroup=res_30m \
  input=L8_2015_01,L8_2015_02,L8_2015_03,L8_2015_04,L8_2015_05,L8_2015_06,L8_2015_07
i.cluster group=L8_2015 subgroup=res_30m \
  signaturefile=cluster_L8_2015 \
  classes=10 reportfile=rep_clust_L8_2015.txt --overwrite
i.maxlik group=L8_2015 subgroup=res_30m \
  signaturefile=cluster_L8_2015 \
  output=L8_2015_cluster_classes reject=L8_2015_cluster_reject --overwrite
\end{lstlisting}
\end{listing}\vspace{-9pt}

Mapping and visualization of the output data was performed using a combination of the GRASS GIS modules `d.rast', `d.legend', `r.colors', `d.mon', `g.region', and `d.out.file', as~shown for the Landsat image 2023 for a map produced using {the }classification and accuracy assessment {with} rejection probability classes (Listing \ref{lst04}). The~same procedure was repeated for images taken in 2018 and 2023{, respectively}.

\begin{listing}[H]
\caption{GRASS GIS code for grouping bands and clustering.} 
\begin{lstlisting}[language=bash,caption={},style=mystyle,label={lst04}]
d.mon wx0
g.region raster=L9_2023_cluster_classes -p
r.colors L9_2023_cluster_classes color=roygbiv -e
d.rast L9_2023_cluster_classes
d.legend raster=L9_2023_cluster_classes title=``24 August 2023'' title_fontsize=14 font=``Helvetica'' fontsize=12 bgcolor=white border_color=white
d.out.file output=Mozambique_2023 format=jpg --overwrite
\end{lstlisting}
\end{listing}\vspace{-12pt}


{The }accuracy assessment for the unsupervised classification was performed using the 'reject' function of the `i.maxlik' modules of GRASS GIS, which calculates the output raster map holding {the} reject threshold results. The~aim of this step is to evaluate the confidence level at which each cell in the raster image is categorized during the classification process. {{The} algorithm} generates the reject threshold map layer, which contains the index to a computed confidence level for each classified pixel in the{ }satellite scene. The~{predefined confidence} intervals {include} sixteen {values}. The~rejection probability image indicated the values of pixels between the acceptable at one, i.e.,~the pixel is kept up to 16, which means that the pixel is~rejected.

The final part of the workflow included the machine learning (ML) approach of random forest (RF) classification of GRASS GIS, which was performed using a set of modules{: }`r.random', `r.learn.train', `r.learn.predict', `r.category', and auxiliary modules of data processing and visualization ({ }Listing \ref{lst05}). 

\begin{listing}[H]
\caption{GRASS GIS code for random forest classification of the Landsat 8--9 OLI/TIRS images using machine learning (ML) approach of remote sensing data processing.} 
\begin{lstlisting}[language=bash,caption={},style=mystyle,label={lst05}]
g.list rast
g.region raster=L8_2015_01 -p
r.random input=L8_2015_cluster_classes seed=100 npoints=1000 raster=training_pixels
# Next, create the imagery group with all Landsat-8 OLI/TIRS bands:
i.group group=L8_2015 input=L8_2015_01,L8_2015_02,L8_2015_03,L8_2015_04,L8_2015_05,L8_2015_06,L8_2015_07 --overwrite
r.learn.train group=L8_2015 training_map=training_pixels model_name=RandomForestClassifier n_estimators=500 save_model=rf_model.gz --overwrite
r.learn.predict group=L8_2015 load_model=rf_model.gz output=rf_classification
# check raster categories - they are automatically applied to the classification output
r.category rf_classification
# copy color scheme from landclass training map to result and display
r.colors rf_classification raster=L8_2015_classes_roi
d.mon wx1
r.colors rf_classification color=plasma -e
d.rast rf_classification
d.legend raster=rf_classification title=``Random Forest: 2015'' title_fontsize=14 font=``Helvetica'' fontsize=12 bgcolor=white border_color=white
d.out.file output=RF_Mozambique_2015 format=jpg --overwrite
\end{lstlisting}
\end{listing}\vspace{-19pt}

The~training of a random forest classification model was performed using `r.learn.train'. The~training {data} were generated from {the {previous }step of}{ }land cover classification{. {The} scheme of} {the }training pixels was applied to perform {the} classification on the target {satellite }image. The~prediction was performed using the `r.learn.predict' module {of the GRASS GIS}. 

%----------- section ----------->
\section{Results}

The color composites created using a combination of bands for false and natural color composites are shown in Figure~\ref{fig05} with the  case of an image from 2015. The~image dates were distributed throughout the important stages of the vegetation growth period in~order to obtain temporal profiles of each kind of land cover type {to analyze} the separability among classes. {The particular characteristics of Mozambique consists of its longitudinal extent, with {the} coastline stretching for ca. 2000 km{, }11\degree~S to 27\degree~S. Moreover, the~effects from tropical ocean currents were {directed} southwards along the length of the country {influence} environmental setting. Therefore, regional characteristics naturally differ within the country. 

\begin{figure}[H]
\hspace{15pt}
\captionsetup[subfigure]{justification=centering}
\makebox[0.90\linewidth][r]{%
	\begin{subfigure}[b]{.40\textwidth}
		\centering
			\includegraphics[width=0.87\textwidth]{Fig_05a.jpg}
		\caption{2015}
	\end{subfigure}%
	\begin{subfigure}[b]{.40\textwidth}
		\centering
		\includegraphics[width=0.87\textwidth]{Fig_05b.jpg}
		\caption{2018}
	\end{subfigure}%
	\begin{subfigure}[b]{.40\textwidth}
		\centering
		\includegraphics[width=0.87\textwidth]{Fig_05c.jpg}
		\caption{2023}
	\end{subfigure}%
}
\caption{Color composites of the Landsat-8 OLI/TIRS image on 2015 on the coastal region of Beira, Mozambique: (\textbf{a}) Bands 7-5-3; (\textbf{b}) Bands 2-3-4; (\textbf{c}) Bands 3-4-5. Major band designations for the Landsat-8 are Band 1---coastal aerosol; Band 2---blue;  Band 3---green; Band 4---red; Band 5---near infrared (NIR); Band 6---shortwave infrared (SWIR) 1; Band 7---shortwave infrared (SWIR) 2.}\label{fig05}
\end{figure}

\vspace{-6pt}

{Weather} patterns {in} Mozambique {vary} due to the global warming and increase in temperature. Nevertheless,{ }the~whole country broadly follows a southern African weather pattern{:}~{ }wet period with heavy rainy between November and March and dry period{ }from April to late November. }Since Mozambique has a tropical climate which has two distinct seasons, the~images were taken during the dry period. {Better} cloudiness conditions {and} dense vegetation~coverage were observed{ during }dry season,{ as shown in} the~images{.}  

The results of the classification {are} shown in Figure~\ref{fig06}{. They }were used as training data for the next step of random forest (RF){ modeling, }{which requires} seed data. The~land cover types{ }include {the following classes: }(1) mosaic cropland vegetation, (2) artificial surfaces and associated areas; (3) rain-fed croplands; (4) open grasslands; (5) sparse vegetation; (6) broad-leaved deciduous woodland; (7) mosaic forests, grassland and shrubland; (8) water areas; (9) needle-leaved deciduous evergreen forests; and (10) repeatedly or permanently flooded lands in brackish waters with broad-leaved forests or shrub lands. {The }coastal areas are characterized by{ }grassland or woody vegetation distributed in {areas of }{ }soil regularly flooded {by} fresh, brackish, or saline~water. 

Figure~\ref{fig07} shows the result of a chi square test on the accuracy {assessment} of the classification on each discriminant result at various threshold levels of confidence. {The aim of the chi square test is to determine if a correlation exists between two qualitative variables {and} if the existing associations are statistically significant. }When using the maximum likelihood classification method to classify the land cover types, the~{accuracy of the classification} depends on the quality of the training dataset and technical parameters of the {satellite }image. Thus, evaluating the accuracy analysis{ }{presents an }important {step in {classification}}. The~approach to accuracy assessments{ }is based on the function embedded in `i.maxlik', where{ }pixels are classified according to the probability of their correct assignment to {each} land cover~class{.} 

\begin{figure}[H]
\hspace{15pt}
\captionsetup[subfigure]{justification=centering}
\makebox[0.90\linewidth][r]{%
	\begin{subfigure}[b]{.40\textwidth}
		\centering
			\includegraphics[width=0.97\textwidth]{Fig_06a.jpg}
		\caption{2015}
	\end{subfigure}%
	\begin{subfigure}[b]{.40\textwidth}
		\centering
		\includegraphics[width=0.97\textwidth]{Fig_06b.jpg}
		\caption{2018}
	\end{subfigure}%
	\begin{subfigure}[b]{.40\textwidth}
		\centering
		\includegraphics[width=0.97\textwidth]{Fig_06c.jpg}
		\caption{2023}
	\end{subfigure}%
}
\caption{Classified %MDPI: Please add the left bracket in the image, e.g., "1)" should be "(1)". % Author's reply: corrected. Figure 6 is updated as requested.
 Landsat-8--9 OLI/TIRS images on the coastal region of Beira, Mozambique, partitioned into 10 classes using k-means clustering approach: (\textbf{a}) 2015; (\textbf{b}) 2018; (\textbf{c}) 2023.}\label{fig06}
\end{figure}

{As}~can be seen in Figure~\ref{fig07} {showing} reject threshold results, {the }lowest rejection values {on }the~raster maps for three years---2015, 2018 and 2023---indicate {the correctly} classified pixels{ }while the highest rejection values indicate the erroneously classified pixels{.} The~analysis of Figure~\ref{fig07} indicates that the {majority of the} correctly indicated pixels {belongs} to the river and inner water classes as well as vegetation types with distinct separability properties{. In~contrast,} the least secure pixels belong to the tidal region in the coastal waters due to the spectral properties of shelf water areas where pixels {might} be misclassified {and }with other land cover types {due to the high similarity of spectral reflectance values}. 

\begin{figure}[H]
\hspace{15pt}
\captionsetup[subfigure]{justification=centering}
\makebox[0.90\linewidth][r]{%
	\begin{subfigure}[b]{.40\textwidth}
		\centering
			\includegraphics[width=0.97\textwidth]{Fig_07a.jpg}
		\caption{2015}
	\end{subfigure}%
	\begin{subfigure}[b]{.40\textwidth}
		\centering
		\includegraphics[width=0.97\textwidth]{Fig_07b.jpg}
		\caption{2018}
	\end{subfigure}%
	\begin{subfigure}[b]{.40\textwidth}
		\centering
		\includegraphics[width=0.97\textwidth]{Fig_07c.jpg}
		\caption{2023}
	\end{subfigure}%
}
\caption{Accuracy %MDPI: Please add the left bracket in the image, e.g., "1)" should be "(1)". % Author's reply: corrected. Figure 7 is updated as requested.
 assessment using the rejection probability for pixels classified using k-means clustering method for Landsat 8--9 OLI/TIRS images on the coastal region of Beira, Mozambique: (\textbf{a})~2015; (\textbf{b}) 2018; (\textbf{c}) 2023.}\label{fig07}
\end{figure}

\vspace{-6pt}

This method is suitable for providing {a} sufficient estimation on the accuracy of training samples that were used for the final steps of {classification using random forest}{. }The resulting images of the random forest classification of Landsat {8--9} OLI/TIRS images by the ML approach of GRASS GIS for the years 2015, 2018, and 2023 are presented in Figure~\ref{fig08}. 

\begin{figure}[H]
\hspace{15pt}
\captionsetup[subfigure]{justification=centering}
\makebox[0.90\linewidth][r]{%
	\begin{subfigure}[b]{.40\textwidth}
		\centering
			\includegraphics[width=0.97\textwidth]{Fig_08a.jpg}
		\caption{2015}
	\end{subfigure}%
	\begin{subfigure}[b]{.40\textwidth}
		\centering
		\includegraphics[width=0.97\textwidth]{Fig_08b.jpg}
		\caption{2018}
	\end{subfigure}%
	\begin{subfigure}[b]{.40\textwidth}
		\centering
		\includegraphics[width=0.97\textwidth]{Fig_08c.jpg}
		\caption{2023}
	\end{subfigure}%
}
\caption{Random %MDPI: Please add the left bracket in the image, e.g., "1)" should be "(1)". % Author's reply: corrected. Figure 8 is updated as requested.
 forest classification using machine learning (ML) approach of GRASS GIS for the Landsat 8--9 OLI/TIRS images on coastal region of Beira, Mozambique: (\textbf{a}) 2015; (\textbf{b}) 2018; (\textbf{c}) 2023.}\label{fig08}
\end{figure}

Figure~\ref{fig08} shows the {results of the land cover type }classification of the coastal region of Mozambique {performed }using the RF classifier approach of the ML method {in} GRASS GIS. The~land cover types in Mozambique were selected for random forest classification with the estimator settings tab, providing access to the most pertinent parameters{. These} affect {the} {previously} described algorithms, where{ }landscape patches were classified using {the{ }}maximum likelihood {classifier}{.} {The}~embedded Scikit-Learn library of Python contents {provide} %Please check meaning
{the} estimator parameters that are supplied for classification. {These} parameters were tuned using a grid search in the Landsat raster maps by inputting multiple regressors{. This enabled us }to indicate the representativeness of {the }land cover classes, uniformity or heterogeneity of patches, and the classification of pixels located far from the edges of the raster~scene. 

{The}{ }ML approach{ }by the Scikit-Learn library of Python {ensures} that the {confusion of} image {pixels} is avoided as much as possible compared to  traditional classification methods. Moreover,{ }training pixels generated during clustering and used to perform {the} RF-based classification separated the images into {clusters} of pixels. These groups were used both as {a} decision tree training seed dataset for random forest classifiers and~as {a} post-classification accuracy assessment. In~this way, the~RF method avoided biased reference information of the group of~pixels.

%----------- section ----------->
\section{Discussion}

In this study, the~ML method of random forest classification of GRASS GIS was applied for {the processing and analysis of }three Landsat 8--9 OLI/TIRS images {with the aim of }land cover mapping {in} the coastal region of the Bight of Sofala, Mozambique. Quantitative image processing and interpretation are not easy {tasks} due to the complex function of {multiple factors. These include, for~instance, }technical parameters of the image (sun, sensor, and target geometry), regional environmental effects (cloudiness, wetness, {heterogeneity, }diversification and fragmentation of landscapes,{ }topographic ruggedness, and slope curvature) {as well as} sensitivity to leaf characteristics~\cite{Mitchard}. The~ML approach of GRASS GIS enabled{ us }to perform a quantitative image classification using the random forest algorithm,{ }which evaluated the similarity of pixels and {assigned} them into classes using {the} `r.learn.train' module. {The} statistical information {was obtained }from the raster data analysis {to }represent{ }the land cover types in Mozambique for{ the }years 2015, 2018, and~2023. 

{Such} an algorithm enabled us to model{ }clusters of pixels{, }which {were} used to compare the climate effects {that cause} changes in land cover classes and {lead to spatio-temporal }landscape {dynamics}. This was achieved using the ensemble classification {achieved by the random forest approach}{. }Each tree in {this} ensemble was based on a random subsample of the training data that were initially obtained using {the }unsupervised classification as a part of data processing. This work is comparable with investigations performed by {the authors of }related works{, for~instance,} those in~\cite{FERNANDO2023,KOLARIK2024111445}. 

{Similar research} ~\cite{BERA20234691,land12122149,ZHANG2023108250} investigated the complementarities of the ML approach for mapping coastal regions using a set of multi-resolution imagery from the Sentinel-2 Multispectral Instrument (MSI) and Landsat-8 OLI. To~continue these studies, the experiments on image processing using the random forest approach reported in this article  investigated the landscape dynamics in the study area of eastern Mozambique{. This was achieved }using a randomly selected subset of the predictors available during each node split in the ML-based classification method of GRASS GIS. Other studies~\cite{MOUNTRAKIS2023106,SINGH2021100645,ZHENG2023100974} used a neural network {as a technical approach to the classification of} vegetation and land cover {types. Such a }technique {of} remote sensing data processing {{aims} to classify the }land cover types using ML methods and {reports} high accuracy. {Hence, {we used the} emerging technologies of scripts and programming algorithms for image processing through a {set} of {GRASS GIS} modules{.}}

{The demonstrated} workflow {was} developed using a short series of {the }Landsat images. These images had various combinations of pixels which were evaluated by the ML{ }methods to analyze{ }the land cover {changes }using prediction analyses. {Each} tree of the algorithm produced a prediction{. The} final result of the image analysis {was} obtained by selecting the mean values across all of the trees in the Landsat image raster using computer vision. In~other studies, a~decision tree classifier was adopted to {the} RS data processing for automatic classification using diverse software. For~instance, ref.~\cite{BROWNDECOLSTOUN2003316} used the ensemble techniques of boosting and consensus filtering of the training data to improve the quality of the training dataset. Ref.~\cite{VAIDYA2024101770} compared the integrated use of the three classification schemes based on RF, DL neural network, and convolutional neural network to distinguish urban, natural, and heterogeneous landscape patches. Likewise{,} the~ExtraTreesClassifier has been used as a{ }variant of the GRASS GIS ML algorithm. Here, each node split during{ }image processing includes the selected splitting rule, which is based on the most optimal~solution. 

Such optimization {applied to} several randomly generated thresholds by the random forest method {presented} advantages over the traditional unsupervised methods of image classification when generating classes for image classification. For~instance, ref.~\cite{PASQUARELLA2023103561} reports that {the} parameter-adjusted decision tree ensembles and {RF} ensembles performed well for image segmentation and supervised and outperformed other approaches. Likewise, the~{updates} in the algorithms were reported using {the }improved decision tree algorithm based on{ }the fuzzy approach{ }for a technical improvement in the{ }performance~\cite{ALOBEIDAT20151192}. Although~the study presented here is not intended to suggest that scripting methods replace the traditional GIS, the~integration of the scripting workflow of GRASS GIS with a satellite image analysis{ }and cartographic plotting provides a more effective approach. {This} study demonstrated that the programming codes{ }supplement the traditional cartographic methods used in the remote sensing software for monitoring coastal areas. Moreover, this work illustrated how the analysis of land cover changes can be visualized in the automatic way using ML algorithms of GRASS GIS. 
%----------- section ----------->
\section{Conclusions}

{{The application of }programming for {RS} data processing {provides} new perspectives {in} cartographic data processing. {Such} approaches are useful for mapping complex landscapes in the entanglement of human-based factors and environmental processes that affect coastal ecosystems. In~view of this, the current study presented advanced mapping approaches to spatial analysis and {visualization} of {the RS} data}. {The} use of the ML and intelligence classifiers of GRASS GIS increased the accuracy and precision of the produced thematic maps through the high level of automation. {Moreover, it} improved the performance of tasks using diversified modules and the application of RF algorithms for {land cover} classifications. {The}~{satellite }images were classified using {the }unsupervised and supervised{ }methods of the ML modules of GRASS GIS, including the random forest~classifier. 

The Bight of Sofala, located in the eastern coastal part of Mozambique, was {selected} as the test area {due to the variability {and} high heterogeneity of landscapes with a complex mosaic of landscape patches}. Ten types of land cover classes were identified on the {Landsat satellite }images {collected in the years 2015, 2018, and 2023. The~data were }classified using {the} random forest algorithm of the machine learning (ML) methods supported by the Scikit-Learn library of Python integrated with GRASS GIS. The~clustering approach of the machine-based classification was {employed} for training and validation of the result of {the} RF classification. The~{benefits} of the{ }{RS} data{ used }for environmental monitoring {consist of the fact} that {they} provide an efficient way to map and visualize remotely located regions{ }such as coastal regions in East Africa. Scripts{ }used as a{ }methodological benchmark of {the} automated image processing and cartographic workflow enabled us to focus on the analysis of the coastal landscapes of Mozambique{.} 

The coastal region of the Bight of Sofala and {its }surroundings were analyzed {using} the classified series of {satellite }images in {Landsat with the applied methods of }automated data analysis of the ML approach. The~{advantage of the }proposed algorithms of GRASS GIS scripts {consists of their repeatability and applicability. }{The use of the }supervised classification by the `r.learn.train' module can be{ }extended to other regions and environmental applications where satellite image processing is required{. In~this way, the {ML} method} classification of landscapes{{ }can be applied in{ }studies that use} a multi-temporal set of{ }images. The~advanced approach of the GRASS GIS scripts was introduced to evaluate the environmental properties of the coastal landscapes of Mozambique{. }Changes in land cover types caused by climate effects in East Africa were visualized using image processing by the GRASS GIS modules  `r.learn.train' , `i.cluster', `i.maxlik', and other auxiliary~modules. 

The{ }results {visualize} the {dynamical }behavior of land cover types in the coastal landscapes of Mozambique{ caused by} climate effects, {which act} as {drivers} for desertification and land degradation in{ }the coastal ecosystems {of East Africa}. {We} also discussed the ways in which the ML methods {technically }contrast with {the }traditional methods of {the }unsupervised classification{. For~instance, these use} k-means clustering {as an approach} to image analysis{ }for evaluating landscape changes using a series of images taken on different dates{. }Thus, {it has been demonstrated that }the ML approach of GRASS GIS{ has }high automation in the sequential image processing workflow, which enables {the }effective use of RS data for environmental monitoring of coastal~regions.

\funding{This publication was funded by the Editorial Office of \emph{Coasts}, Multidisciplinary Digital Publishing Institute (MDPI), by~providing a 100\% discount on the APC of this manuscript.}

\institutionalreview{Not applicable.}

\informedconsent{Not applicable.}

\dataavailability{{The data are contained within the article.} %MDPI: "Not applicable" is only used for review papers or articles for which no new data were created. Normally, the DAS for research articles can state "Data are contained within the article." or "Data are contained within the article and supplementary materials." Please refer to the complete guideline at  https://www.mdpi.com/ethics#_bookmark21. % Author's reply: corrected as recommended.
} 

\acknowledgments{The authors thank the reviewers for reading and reviewing this~manuscript.}

\conflictsofinterest{The author declares no conflicts of~interest.} 

\abbreviations{Abbreviations}{
The following abbreviations are used in this manuscript:\\

\noindent 
\begin{tabular}{@{}ll}
AVHRR & Advanced Very High Resolution Radiometer \\
CCTools & Cooperative Computing Tools \\
DCW & Digital Chart of the World \\
DL & Deep Learning \\
DN & Digital Numbers \\
FAO & Food and Agriculture Organization \\
GDAL & Geospatial Data Abstraction Library \\
GEBCO & General Bathymetric Chart of the Oceans \\ 
GRASS & Geographic Resources Analysis Support System \\
GIS & Geographic Information System \\
GCC & GNU Compiler Collection \\
GMT & Generic Mapping Tools \\
Landsat OLI/TIRS & Landsat Operational Land Imager and Thermal Infrared Sensor \\
Landsat MSS & Multispectral Scanner \\
Landsat TM & Landsat Thematic Mapper\\
ML & Machine Learning \\
NIR & Near Infrared \\
NOAA & National Oceanic and Atmospheric Administration \\
RF & Random Forest \\
Sentinel MSI & Sentinel Multispectral Instrument \\
SPOT & Satellite Pour l'Observation de la Terre \\
SRTM & Shuttle Radar Topography Mission \\
SWIR & Shortwave Infrared \\
ToAR & Top-of-Atmosphere Radiances \\ 
USGS & United States Geological Survey\\
UTM & Universal Transverse Mercator \\
\end{tabular}
}

%%%%%%%%%%%%%%%%%%%%%%%%%%%%%%%%%%%%%%%%%%
\begin{adjustwidth}{-\extralength}{0cm}
%\printendnotes[custom] % Un-comment to print a list of endnotes

\reftitle{References}
%\nocite{*}

%=====================================
% References, variant A: external bibliography
%=====================================
%\bibliography{Bibliography.bib}


\begin{thebibliography}{999}

\bibitem[Fagan and DeFries(2024)]{FAGAN2024432}
Fagan, M.E.; DeFries, R.S.
\newblock Remote Sensing and Image Processing. In {\em Encyclopedia of
  Biodiversity}, 3rd  ed.;  Scheiner, S.M., Ed.;
  Academic Press: Oxford, UK, 2024; pp. 432--445.
\newblock {\url{https://doi.org/10.1016/B978-0-12-822562-2.00060-8}}.

\bibitem[Loeb(2023)]{LOEB2023}
Loeb, N.G.
\newblock Satellites and Satellite Remote Sensing|Earth's Radiation Budget.
  In {\em Reference Module in Earth Systems and Environmental Sciences};
  Elsevier: Amsterdam, The Netherlands, % MDPI:  Newly added information, please confirm.  The following highlights are the same. %<-- Author's reply: yes, confirm.
  2023.
\newblock {\url{https://doi.org/10.1016/B978-0-323-96026-7.00002-3}}.

\bibitem[Brenning(2012)]{6352393}
Brenning, A.
\newblock Spatial cross-validation and bootstrap for the assessment of
  prediction rules in remote sensing: The R package sperrorest.
\newblock In Proceedings of the 2012 IEEE International Geoscience and Remote
  Sensing Symposium, Munich, Germany, 22--27 July 2012; pp. 5372--5375. %<-- Author's reply: yes, confirm.
\newblock {\url{https://doi.org/10.1109/IGARSS.2012.6352393}}.

\bibitem[Lemenkova and Debeir(2023{\natexlab{a}})]{info14040249}
Lemenkova, P.; Debeir, O.
\newblock Recognizing the Wadi Fluvial Structure and Stream Network in the Qena
  Bend of the Nile River, Egypt, on Landsat 8-9 OLI Images.
\newblock {\em Information} {\bf 2023}, {\em 14}.
\newblock {\url{https://doi.org/10.3390/info14040249}}.

\bibitem[Lemenkova and Debeir(2023{\natexlab{b}})]{Lemenkova202324}
Lemenkova, P.; Debeir, O.
\newblock {Time Series Analysis of Landsat Images for Monitoring Flooded Areas
  in the Inner Niger Delta, Mali}.
\newblock {\em Artif. Satell.} {\bf 2023}, {\em 58},~278--313.
\newblock {\url{https://doi.org/10.2478/arsa-2023-0011}}.

\bibitem[Zhao et~al.(2024)Zhao, Huang, Li, Yu, and Shen]{ZHAO2024169152}
Zhao, D.; Huang, J.; Li, Z.; Yu, G.; Shen, H.
\newblock Dynamic monitoring and analysis of chlorophyll-a concentrations in
  global lakes using Sentinel-2 images in Google Earth Engine.
\newblock {\em Sci. Total Environ.} {\bf 2024}, {\em 912},~169152.
\newblock {\url{https://doi.org/10.1016/j.scitotenv.2023.169152}}.

\bibitem[Lemenkova(2020)]{Lemenkova202021}
Lemenkova, P.
\newblock {Sentinel-2 for High Resolution Mapping of Slope-Based Vegetation
  Indices Using Machine Learning by SAGA GIS}.
\newblock {\em Transylv. Rev. Syst. Ecol. Res.}
  {\bf 2020}, {\em 22},~17--34.
\newblock {\url{https://doi.org/10.2478/trser-2020-0015}}.

\bibitem[Choudhary et~al.(2022)Choudhary, Shi, Dong, and
  Paringer]{CHOUDHARY20222443}
Choudhary, K.; Shi, W.; Dong, Y.; Paringer, R.
\newblock Random Forest for rice yield mapping and prediction using Sentinel-2
  data with Google Earth Engine.
\newblock {\em Adv. Space Res.} {\bf 2022}, {\em 70},~2443--2457.
\newblock {\url{https://doi.org/10.1016/j.asr.2022.06.073}}.

\bibitem[Ji and Brown(2017)]{JI2017215}
Ji, L.; Brown, J.F.
\newblock Effect of NOAA satellite orbital drift on AVHRR-derived phenological
  metrics.
\newblock {\em Int. J. Appl. Earth Obs. Geoinf.} {\bf 2017}, {\em 62},~215--223.
\newblock {\url{https://doi.org/10.1016/j.jag.2017.06.013}}.

\bibitem[Zhan and Liang(2022)]{ZHAN2022112836}
Zhan, C.; Liang, S.
\newblock Improved estimation of the global top-of-atmosphere albedo from AVHRR
  data.
\newblock {\em Remote Sens. Environ.} {\bf 2022}, {\em 269},~112836.
\newblock {\url{https://doi.org/10.1016/j.rse.2021.112836}}.

\bibitem[Fensholt et~al.(2009)Fensholt, Rasmussen, Nielsen, and
  Mbow]{FENSHOLT20091886}
Fensholt, R.; Rasmussen, K.; Nielsen, T.T.; Mbow, C.
 Evaluation of earth observation based long term vegetation \mbox{trends---Intercomparing} NDVI time series trend analysis consistency of Sahel from
  AVHRR GIMMS, Terra MODIS and SPOT VGT data.
 {\em Remote Sens. Environ.} {\bf 2009}, {\em
  113},~1886--1898.
\newblock {\url{https://doi.org/10.1016/j.rse.2009.04.004}}.

\bibitem[Chen et~al.(2023)Chen, Yang, Jiang, and Liu]{CHEN2023e19654}
Chen, C.; Yang, X.; Jiang, S.; Liu, Z.
\newblock Mapping and spatiotemporal dynamics of land-use and land-cover change
  based on the Google Earth Engine cloud platform from Landsat imagery: A case
  study of Zhoushan Island, China.
\newblock {\em Heliyon} {\bf 2023}, {\em 9},~e19654.
\newblock {\url{https://doi.org/10.1016/j.heliyon.2023.e19654}}.

\bibitem[Lemenkova(2023)]{Lemenkova202323}
Lemenkova, P.
\newblock {Using open-source software GRASS GIS for analysis of the
  environmental patterns in Lake Chad, Central Africa}.
\newblock {\em Die Bodenkult. J. Land Manag. Food Environ.} {\bf 2023}, {\em 74},~49--64.
\newblock {\url{https://doi.org/10.2478/boku-2023-0005}}.

\bibitem[Lemesios and Petropoulos(2024)]{LEMESIOS2024101153}
Lemesios, I.; Petropoulos, G.P.
\newblock Vegetation regeneration dynamics of a natural mediterranean ecosystem
  following a wildfire exploiting the LANDSAT archive, google earth engine and
  geospatial analysis techniques.
\newblock {\em Remote. Sens. Appl. Soc. Environ.} {\bf
  2024}, \emph{34}, 101153.
\newblock {\url{https://doi.org/10.1016/j.rsase.2024.101153}}.

\bibitem[Senay et~al.(2022)Senay, Friedrichs, Morton, Parrish, Schauer, Khand,
  Kagone, Boiko, and Huntington]{SENAY2022113011}
Senay, G.B.; Friedrichs, M.; Morton, C.; Parrish, G.E.; Schauer, M.; Khand, K.;
  Kagone, S.; Boiko, O.; Huntington, J.
\newblock Mapping actual evapotranspiration using Landsat for the conterminous
  United States: Google Earth Engine implementation and assessment of the
  SSEBop model.
\newblock {\em Remote Sens. Environ.} {\bf 2022}, {\em 275},~113011.
\newblock {\url{https://doi.org/10.1016/j.rse.2022.113011}}.

\bibitem[Lemenkova and Debeir(2023)]{Lemenkova202321}
Lemenkova, P.; Debeir, O.
\newblock {Environmental mapping of Burkina Faso using TerraClimate data and
  satellite images by GMT and R scripts}.
\newblock {\em Adv. Geod. Geoinf.} {\bf 2023}, {\em
  72},~e45.
\newblock {\url{https://doi.org/10.24425/agg.2023.146157}}.

\bibitem[Xie et~al.(2023)Xie, Li, Wulan, Zheng, and Shen]{XIE2023}
Xie, Y.; Li, J.; Wulan, T.; Zheng, Y.; Shen, Z.
\newblock Scale dependence of forest fragmentation and its climate sensitivity
  in a semi-arid mountain: Comparing Landsat, Sentinel and Google Earth data.
\newblock {\em Geogr. Sustain.} {\bf 2023}, \emph{in press}.
\newblock {\url{https://doi.org/10.1016/j.geosus.2023.11.008}}.

\bibitem[Radeloff et~al.(2024)Radeloff, Roy, Wulder, Anderson, Cook, Crawford,
  Friedl, Gao, Gorelick, Hansen, Healey, Hostert, Hulley, Huntington, Johnson,
  Neigh, Lyapustin, Lymburner, Pahlevan, Pekel, Scambos, Schaaf, Strobl,
  Woodcock, Zhang, and Zhu]{RADELOFF2024113918}
Radeloff, V.C.; Roy, D.P.; Wulder, M.A.; Anderson, M.; Cook, B.; Crawford,
  C.J.; Friedl, M.; Gao, F.; Gorelick, N.; Hansen, M.;  et~al.
\newblock Need and vision for global medium-resolution Landsat and Sentinel-2
  data products.
\newblock {\em Remote Sens. Environ.} {\bf 2024}, {\em 300},~113918.
\newblock {\url{https://doi.org/10.1016/j.rse.2023.113918}}.

\bibitem[Giri et~al.(2005)Giri, Zhu, and Reed]{GIRI2005123}
Giri, C.; Zhu, Z.; Reed, B.
\newblock A comparative analysis of the Global Land Cover 2000 and MODIS land
  cover data sets.
\newblock {\em Remote Sens. Environ.} {\bf 2005}, {\em 94},~123--132.
\newblock {\url{https://doi.org/10.1016/j.rse.2004.09.005}}.

\bibitem[Wang et~al.(2023)Wang, Sun, Cao, Wang, Zhang, and Cheng]{WANG2023311}
Wang, Y.; Sun, Y.; Cao, X.; Wang, Y.; Zhang, W.; Cheng, X.
\newblock A review of regional and Global scale Land Use/Land Cover (LULC)
  mapping products generated from satellite remote sensing.
\newblock {\em ISPRS J. Photogramm. Remote Sens.} {\bf 2023},
  {\em 206},~311--334.
\newblock {\url{https://doi.org/10.1016/j.isprsjprs.2023.11.014}}.

\bibitem[He et~al.(2018)He, Liang, Wang, Cao, Gao, Yu, and Feng]{HE2018181}
He, T.; Liang, S.; Wang, D.; Cao, Y.; Gao, F.; Yu, Y.; Feng, M.
\newblock Evaluating land surface albedo estimation from Landsat MSS, TM, ETM+,
  and OLI data based on the unified direct estimation approach.
\newblock {\em Remote Sens. Environ.} {\bf 2018}, {\em 204},~181--196.
\newblock {\url{https://doi.org/10.1016/j.rse.2017.10.031}}.

\bibitem[Hussain et~al.(2022)Hussain, Mubeen, and
  Karuppannan]{HUSSAIN2022103117}
Hussain, S.; Mubeen, M.; Karuppannan, S.
\newblock Land use and land cover (LULC) change analysis using TM, ETM+ and OLI
  Landsat images in district of Okara, Punjab, Pakistan.
\newblock {\em Phys. Chem. Earth Parts A/B/C} {\bf 2022},
  {\em 126},~103117.
\newblock {\url{https://doi.org/10.1016/j.pce.2022.103117}}.

\bibitem[Estoque and Murayama(2015)]{ESTOQUE2015205}
Estoque, R.C.; Murayama, Y.
\newblock Classification and change detection of built-up lands from Landsat-7
  ETM+ and Landsat-8 OLI/TIRS imageries: A comparative assessment of various
  spectral indices.
\newblock {\em Ecol. Indic.} {\bf 2015}, {\em 56},~205--217.
\newblock {\url{https://doi.org/10.1016/j.ecolind.2015.03.037}}.

\bibitem[Wulder et~al.(2022)Wulder, Roy, Radeloff, Loveland, Anderson, Johnson,
  Healey, Zhu, Scambos, Pahlevan, Hansen, Gorelick, Crawford, Masek,
  Hermosilla, White, Belward, Schaaf, Woodcock, Huntington, Lymburner, Hostert,
  Gao, Lyapustin, Pekel, Strobl, and Cook]{WULDER2022113195}
Wulder, M.A.; Roy, D.P.; Radeloff, V.C.; Loveland, T.R.; Anderson, M.C.;
  Johnson, D.M.; Healey, S.; Zhu, Z.; Scambos, T.A.; Pahlevan, N.;  et~al.
\newblock Fifty years of Landsat science and impacts.
\newblock {\em Remote Sens. Environ.} {\bf 2022}, {\em 280},~113195.
\newblock {\url{https://doi.org/10.1016/j.rse.2022.113195}}.

\bibitem[Ferreira et~al.(2009)Ferreira, Andrade, Bandeira, Cardoso, Mendes, and
  Paula]{Ferreira}
Ferreira, M.A.; Andrade, F.; Bandeira, S.O.; Cardoso, P.; Mendes, R.N.; Paula,
  J.
\newblock Analysis of cover change (1995–2005) of Tanzania/Mozambique
  trans-boundary mangroves using Landsat imagery.
\newblock {\em Aquat. Conserv. Mar. Freshw. Ecosyst.} {\bf
  2009}, {\em 19},~S38--S45.
\newblock {\url{https://doi.org/10.1002/aqc.1042}}.

\bibitem[Montfort et~al.(2021)Montfort, Bégué, Leroux, Blanc, Gond, Cambule,
  Remane, and Grinand]{Montfort}
Montfort, F.; Bégué, A.; Leroux, L.; Blanc, L.; Gond, V.; Cambule, A.H.;
  Remane, I.A.D.; Grinand, C.
\newblock From land productivity trends to land degradation assessment in
  Mozambique: Effects of climate, human activities and stakeholder definitions.
\newblock {\em Land Degrad. Dev.} {\bf 2021}, {\em 32},~49--65.
\newblock {\url{https://doi.org/10.1002/ldr.3704}}.

\bibitem[Fatoyinbo et~al.(2008)Fatoyinbo, Simard, Washington-Allen, and
  Shugart]{Fatoyinbo}
Fatoyinbo, T.E.; Simard, M.; Washington-Allen, R.A.; Shugart, H.H.
\newblock Landscape-scale extent, height, biomass, and carbon estimation of
  Mozambique's mangrove forests with Landsat ETM+ and Shuttle Radar Topography
  Mission elevation data.
\newblock {\em J. Geophys. Res. Biogeosci.} {\bf 2008},
  {\em 113%MDPI: Please add the page number. % Author's reply: added (1-13)
},~1--13.
\newblock {\url{https://doi.org/10.1029/2007JG000551}}.

\bibitem[Ribeiro et~al.(2008)Ribeiro, Saatchi, Shugart, and
  Washington-Allen]{Ribeiro}
Ribeiro, N.S.; Saatchi, S.S.; Shugart, H.H.; Washington-Allen, R.A.
\newblock Aboveground biomass and leaf area index (LAI) mapping for Niassa
  Reserve, northern Mozambique.
\newblock {\em J. Geophys. Res. Biogeosci.} {\bf 2008},
 {\em 113%MDPI: Please add the page number. % Author's reply: added (1-13)
},~1--12.
\newblock {\url{https://doi.org/10.1029/2007JG000550}}.

\bibitem[Li et~al.(2023)Li, Yigitcanlar, Nepal, Nguyen, and Dur]{LI2023104653}
Li, F.; Yigitcanlar, T.; Nepal, M.; Nguyen, K.; Dur, F.
\newblock Machine learning and remote sensing integration for leveraging urban
  sustainability: A review and framework.
\newblock {\em Sustain. Cities Soc.} {\bf 2023}, {\em 96},~104653.
\newblock {\url{https://doi.org/10.1016/j.scs.2023.104653}}.

\bibitem[Pham et~al.(2023)Pham, Ha, Saintilan, Skidmore, Phan, Le, Viet,
  Takeuchi, and Friess]{PHAM2023104501}
Pham, T.D.; Ha, N.T.; Saintilan, N.; Skidmore, A.; Phan, D.C.; Le, N.N.; Viet,
  H.L.; Takeuchi, W.; Friess, D.A.
\newblock Advances in Earth observation and machine learning for quantifying
  blue carbon.
\newblock {\em Earth-Sci. Rev.} {\bf 2023}, {\em 243},~104501.
\newblock {\url{https://doi.org/10.1016/j.earscirev.2023.104501}}.

\bibitem[Han et~al.(2023)Han, Zhang, Wang, Wang, Huang, Li, Wang, Chen, Li,
  Feng, Fan, Zhang, and Wang]{HAN202387}
Han, W.; Zhang, X.; Wang, Y.; Wang, L.; Huang, X.; Li, J.; Wang, S.; Chen, W.;
  Li, X.; Feng, R.;  et~al.
\newblock A survey of machine learning and deep learning in remote sensing of
  geological environment: Challenges, advances, and opportunities.
\newblock {\em ISPRS J. Photogramm. Remote Sens.} {\bf 2023},
  {\em 202},~87--113.
\newblock {\url{https://doi.org/10.1016/j.isprsjprs.2023.05.032}}.

\bibitem[Khan et~al.(2022)Khan, Vibhute, Mali, and Patil]{KHAN2022101678}
Khan, A.; Vibhute, A.D.; Mali, S.; Patil, C.
\newblock A systematic review on hyperspectral imaging technology with a
  machine and deep learning methodology for agricultural applications.
\newblock {\em Ecol. Informatics} {\bf 2022}, {\em 69},~101678.
\newblock {\url{https://doi.org/10.1016/j.ecoinf.2022.101678}}.

\bibitem[Mullissa et~al.(2023)Mullissa, Reiche, and Herold]{MULLISSA2023113799}
Mullissa, A.; Reiche, J.; Herold, M.
\newblock Deep learning and automatic reference label harvesting for Sentinel-1
  SAR-based rapid tropical dry forest disturbance mapping.
\newblock {\em Remote Sens. Environ.} {\bf 2023}, {\em 298},~113799.
\newblock {\url{https://doi.org/10.1016/j.rse.2023.113799}}.

\bibitem[Lemenkova(2022)]{Lemenkova202206}
Lemenkova, P.
\newblock {Tanzania Craton, Serengeti Plain and Eastern Rift Valley: Mapping of
  geospatial data by scripting techniques}.
\newblock {\em Est. J. Earth Sci.} {\bf 2022}, {\em
  71},~61--79.
\newblock {\url{https://doi.org/10.3176/earth.2022.05}}.

\bibitem[Lemenkova and Debeir(2022)]{Lemenkova202227}
Lemenkova, P.; Debeir, O.
\newblock {Satellite Image Processing by Python and R Using Landsat 9 OLI/TIRS
  and SRTM DEM Data on Côte d’Ivoire, West Africa}.
\newblock {\em J. Imaging} {\bf 2022}, {\em 8},~317.
\newblock {\url{https://doi.org/10.3390/jimaging8120317}}.

\bibitem[Liu et~al.(2023)Liu, Qiu, Feng, Wong, Tsou, Wang, and
  Zhang]{rs15235559}
Liu, J.; Qiu, Z.; Feng, J.; Wong, K.P.; Tsou, J.Y.; Wang, Y.; Zhang, Y.
\newblock Monitoring Total Suspended Solids and Chlorophyll-a Concentrations in
  Turbid Waters: A Case Study of the Pearl River Estuary and Coast Using
  Machine Learning.
\newblock {\em Remote Sens.} {\bf 2023}, {\em 15}, 5559.
\newblock {\url{https://doi.org/10.3390/rs15235559}}.

\bibitem[Wu et~al.(2023)Wu, Jiang, Xu, Yeh, Xu, Wang, and Su]{WU2023103407}
Wu, Y.; Jiang, N.; Xu, Y.; Yeh, T.K.; Xu, T.; Wang, Y.; Su, W.
\newblock Improving the capability of water vapor retrieval from Landsat 8
  using ensemble machine learning.
\newblock {\em Int. J. Appl. Earth Obs. Geoinf.} {\bf 2023}, {\em 122},~103407.
\newblock {\url{https://doi.org/10.1016/j.jag.2023.103407}}.

\bibitem[Chen et~al.(2023)Chen, Huang, Zhang, Chen, and Chen]{CHEN2023101499}
Chen, J.; Huang, J.; Zhang, X.; Chen, J.; Chen, X.
\newblock Monitoring total suspended solids concentration in Poyang Lake via
  machine learning and Landsat images.
\newblock {\em J. Hydrol. Reg. Stud.} {\bf 2023}, {\em
  49},~101499.
\newblock {\url{https://doi.org/10.1016/j.ejrh.2023.101499}}.

\bibitem[Wu and Pan(2023)]{rs15225387}
Wu, Y.; Pan, J.
\newblock Detecting Changes in Impervious Surfaces Using Multi-Sensor Satellite
  Imagery and Machine Learning Methodology in a Metropolitan Area.
\newblock {\em Remote Sens.} {\bf 2023}, {\em 15}, 5387.
\newblock {\url{https://doi.org/10.3390/rs15225387}}.

\bibitem[Lin et~al.(2021)Lin, Li, Yu, Hu, Zhang, Ye, Syed, and
  Li]{LIN2021102370}
Lin, Y.; Li, L.; Yu, J.; Hu, Y.; Zhang, T.; Ye, Z.; Syed, A.; Li, J.
\newblock An optimized machine learning approach to water pollution variation
  monitoring with time-series Landsat images.
\newblock {\em Int. J. Appl. Earth Obs. Geoinf.} {\bf 2021}, {\em 102},~102370.
\newblock {\url{https://doi.org/10.1016/j.jag.2021.102370}}.

\bibitem[Liao(2023)]{LIAO2023110231}
Liao, Q.
\newblock Intelligent classification model of land resource use using deep
  learning in remote sensing images.
\newblock {\em Ecol. Model.} {\bf 2023}, {\em 475},~110231.
\newblock {\url{https://doi.org/10.1016/j.ecolmodel.2022.110231}}.

\bibitem[Guo et~al.(2023)Guo, Xu, Zeng, Liu, and Zhu]{GUO20231}
Guo, J.; Xu, Q.; Zeng, Y.; Liu, Z.; Zhu, X.X.
\newblock Nationwide urban tree canopy mapping and coverage assessment in
  Brazil from high-resolution remote sensing images using deep learning.
\newblock {\em ISPRS J. Photogramm. Remote Sens.} {\bf 2023},
  {\em 198},~1--15.
\newblock {\url{https://doi.org/10.1016/j.isprsjprs.2023.02.007}}.

\bibitem[Zhang et~al.(2023)Zhang, Wang, Wang, Song, Wei, and Hu]{rs15215121}
Zhang, Q.; Wang, G.; Wang, G.; Song, W.; Wei, X.; Hu, Y.
\newblock Identifying Winter Wheat Using Landsat Data Based on Deep Learning
  Algorithms in the North China Plain.
\newblock {\em Remote Sens.} {\bf 2023}, {\em 15}, 5121.
\newblock {\url{https://doi.org/10.3390/rs15215121}}.

\bibitem[Dou et~al.(2021)Dou, Shen, Li, and Guan]{DOU2021102477}
Dou, P.; Shen, H.; Li, Z.; Guan, X.
\newblock Time series remote sensing image classification framework using
  combination of deep learning and multiple classifiers system.
\newblock {\em Int. J. Appl. Earth Obs. Geoinf.} {\bf 2021}, {\em 103},~102477.
\newblock {\url{https://doi.org/10.1016/j.jag.2021.102477}}.

\bibitem[Boulila et~al.(2021)Boulila, Sellami, Driss, Al-Sarem, Safaei, and
  Ghaleb]{BOULILA2021106014}
Boulila, W.; Sellami, M.; Driss, M.; Al-Sarem, M.; Safaei, M.; Ghaleb, F.A.
\newblock RS-DCNN: A novel distributed convolutional-neural-networks
  based-approach for big remote-sensing image classification.
\newblock {\em Comput. Electron. Agric.} {\bf 2021}, {\em
  182},~106014.
\newblock {\url{https://doi.org/10.1016/j.compag.2021.106014}}.

\bibitem[Lin et~al.(2024)Lin, Zhao, Wang, and Tang]{LIN2024102618}
Lin, J.; Zhao, Y.; Wang, S.; Tang, Y.
\newblock A robust training method for object detectors in remote sensing
  image.
\newblock {\em Displays} {\bf 2024}, {\em 81},~102618.
\newblock {\url{https://doi.org/10.1016/j.displa.2023.102618}}.

\bibitem[Roberts et~al.(2022)Roberts, Mwangi, Mukabi, Njui, Nzioka, Ndambiri,
  Bispo, Espirito-Santo, Gou, Johnson, Louis, Pacheco-Pascagaza,
  Rodriguez-Veiga, Tansey, Upton, Robb, and Balzter]{ROBERTS2022105192}
Roberts, J.; Mwangi, R.; Mukabi, F.; Njui, J.; Nzioka, K.; Ndambiri, J.; Bispo,
  P.; Espirito-Santo, F.; Gou, Y.; Johnson, S.;  et~al.
\newblock Pyeo: A Python package for near-real-time forest cover change
  detection from Earth observation using machine learning.
\newblock {\em Comput. Geosci.} {\bf 2022}, {\em 167},~105192.
\newblock {\url{https://doi.org/10.1016/j.cageo.2022.105192}}.

\bibitem[Guo et~al.(2023)Guo, Guo, Zhang, He, Li, Yang, and
  Zhang]{GUO2023e21542}
Guo, Z.; Guo, F.; Zhang, Y.; He, J.; Li, G.; Yang, Y.; Zhang, X.
\newblock A python system for regional landslide susceptibility assessment by
  integrating machine learning models and its application.
\newblock {\em Heliyon} {\bf 2023}, {\em 9},~e21542.
\newblock {\url{https://doi.org/10.1016/j.heliyon.2023.e21542}}.

\bibitem[Redoloza et~al.(2023)Redoloza, Williamson, Headman, and
  Allred]{Redoloza}
Redoloza, F.S.; Williamson, T.N.; Headman, A.O.; Allred, B.J.
\newblock Machine-learning model to delineate sub-surface agricultural drainage
  from satellite imagery.
\newblock {\em J. Environ. Qual.} {\bf 2023}, {\em
  52},~907--921.
\newblock {\url{https://doi.org/10.1002/jeq2.20493}}.

\bibitem[{GRASS Development Team}(2022)]{GRASS_GIS_software}
{GRASS Development Team}.
\newblock {\em Geographic Resources Analysis Support System (GRASS GIS)
  Software}, Version 8.2.;
\newblock Open Source Geospatial Foundation: Beaverton, OR, USA, 2022. % Author's reply: yes, confirm.

\bibitem[Strigaro et~al.(2016)Strigaro, Moretti, Mattavelli, Frigerio, Amicis,
  and Maggi]{STRIGARO201668}
Strigaro, D.; Moretti, M.; Mattavelli, M.; Frigerio, I.; Amicis, M.D.; Maggi,
  V.
\newblock A GRASS GIS module to obtain an estimation of glacier behavior under
  climate change: A pilot study on Italian glacier.
\newblock {\em Comput. Geosci.} {\bf 2016}, {\em 94},~68--76.
\newblock {\url{https://doi.org//10.1016/j.cageo.2016.06.009}}.

\bibitem[Lemenkova(2023)]{Lemenkova202322}
Lemenkova, P.
\newblock {Monitoring Seasonal Fluctuations in Saline Lakes of Tunisia Using
  Earth Observation Data Processed by GRASS GIS}.
\newblock {\em Land} {\bf 2023}, {\em 12},~1995.
\newblock {\url{https://doi.org/10.3390/land12111995}}.

\bibitem[Jasiewicz(2011)]{JASIEWICZ20111525}
Jasiewicz, J.
\newblock A new GRASS GIS fuzzy inference system for massive data analysis.
\newblock {\em Comput. Geosci.} {\bf 2011}, {\em 37},~1525--1531.
\newblock {\url{https://doi.org/10.1016/j.cageo.2010.09.008}}.

\bibitem[Lemenkova(2023)]{Lemenkova202313}
Lemenkova, P.
\newblock {A GRASS GIS Scripting Framework for Monitoring Changes in the
  Ephemeral Salt Lakes of Chotts Melrhir and Merouane, Algeria}.
\newblock {\em Appl. Syst. Innov.} {\bf 2023}, {\em 6},~61.
\newblock {\url{https://doi.org/10.3390/asi6040061}}.

\bibitem[Jasiewicz and Metz(2011)]{JASIEWICZ20111162}
Jasiewicz, J.; Metz, M.
\newblock A new GRASS GIS toolkit for Hortonian analysis of drainage networks.
\newblock {\em Comput. Geosci.} {\bf 2011}, {\em 37},~1162--1173.
\newblock {\url{https://doi.org//10.1016/j.cageo.2011.03.003}}.

\bibitem[Lemenkova(2023)]{Lemenkova202320}
Lemenkova, P.
\newblock {Image Segmentation of the Sudd Wetlands in South Sudan for
  Environmental Analytics by GRASS GIS Scripts}.
\newblock {\em Analytics} {\bf 2023}, {\em 2},~745--780.
\newblock {\url{https://doi.org/10.3390/analytics2030040}}.

\bibitem[Rocchini et~al.(2013)Rocchini, Delucchi, Bacaro, Cavallini, Feilhauer,
  Foody, He, Nagendra, Porta, Ricotta, Schmidtlein, Spano, Wegmann, and
  Neteler]{ROCCHINI201382}
Rocchini, D.; Delucchi, L.; Bacaro, G.; Cavallini, P.; Feilhauer, H.; Foody,
  G.M.; He, K.S.; Nagendra, H.; Porta, C.; Ricotta, C.;  et~al.
\newblock Calculating landscape diversity with information-theory based
  indices: A GRASS GIS solution.
\newblock {\em Ecol. Informatics} {\bf 2013}, {\em 17},~82--93. %MDPI: We removed the redundant issue content, please confirm. % Author's reply: yes, confirm.
  {\url{https://doi.org/10.1016/j.ecoinf.2012.04.002}}.

\bibitem[Hofierka et~al.(2009)Hofierka, Mitášová, and
  Neteler]{HOFIERKA2009387}
Hofierka, J.; Mitášová, H.; Neteler, M.
\newblock Chapter 17 Geomorphometry in GRASS GIS. In {\em Geomorphometry};
  Hengl, T.; Reuter, H.I., Eds.; Developments in
  Soil Science Series; Elsevier: Amsterdam, The Netherlands, 2009; Volume~33, pp. 387--410. % Author's reply: yes, confirm.
\newblock {\url{https://doi.org/10.1016/S0166-2481(08)00017-2}}.

\bibitem[Pedregosa et~al.(2011)Pedregosa, Varoquaux, Gramfort, Michel, and
  Thirion]{Pedregosa}
Pedregosa.; Varoquaux, G.; Gramfort, A.; Michel, V.; Thirion, B.
\newblock {Scikit-learn: Machine Learning in Python}.
\newblock {\em J. Mach. Learn. Res.} {\bf 2011}, {\em
  12},~2825--2830.

\bibitem[Bacar and Faque(2023)]{BACAR2023}
Bacar, F.F.; Faque, H.B.
\newblock Forest holds high rodent diversity than other habitats under a
  rapidly changing and fragmenting landscape in Quirimbas National Park,
  Mozambique.
\newblock {\em Acta Ecol. Sin.} {\bf 2023}.
\newblock {\url{https://doi.org/10.1016/j.chnaes.2023.11.005}}.

\bibitem[Jansen et~al.(2008)Jansen, Bagnoli, and Focacci]{JANSEN2008308}
Jansen, L.J.; Bagnoli, M.; Focacci, M.
\newblock Analysis of land-cover/use change dynamics in Manica Province in
  Mozambique in a period of transition (1990–2004).
\newblock {\em For. Ecol. Manag.} {\bf 2008}, {\em 254},~308--326.
\newblock {\url{https://doi.org/10.1016/j.foreco.2007.08.017}}.

\bibitem[Mucova et~al.(2018)Mucova, Filho, Azeiteiro, and
  Pereira]{MUCOVA2018e00447}
Mucova, S.A.R.; Filho, W.L.; Azeiteiro, U.M.; Pereira, M.J.
\newblock Assessment of land use and land cover changes from 1979 to 2017 and
  biodiversity \& land management approach in Quirimbas National Park, Northern
  Mozambique, Africa.
\newblock {\em Glob. Ecol. Conserv.} {\bf 2018}, {\em 16},~e00447.
\newblock {\url{https://doi.org/10.1016/j.gecco.2018.e00447}}.

\bibitem[Dunham(1990)]{Dunham}
Dunham, K.M.
\newblock Biomass dynamics of herbaceous vegetation in Zambezi riverine
  woodlands.
\newblock {\em Afr. J. Ecol.} {\bf 1990}, {\em 28},~200--212.
\newblock {\url{https://doi.org/10.1111/j.1365-2028.1990.tb01153.x}}.

\bibitem[Dunham(1991)]{Dunham1991}
Dunham, K.M.
\newblock Phenology of Acacia albida trees in Zambezi riverine woodlands.
\newblock {\em Afr. J. Ecol.} {\bf 1991}, {\em 29},~118--129.
\newblock {\url{https://doi.org/10.1111/j.1365-2028.1991.tb00992.x}}.

\bibitem[Simasiku et~al.(2021)Simasiku, Hay, and Abah]{Simasiku}
Simasiku, E.K.; Hay, C.; Abah, J.
\newblock Effect of water level and water quality on small-sized and juvenile
  fish assemblages in the littoral zones of the Zambezi/Chobe floodplain.
\newblock {\em Afr. J. Ecol.} {\bf 2021}, {\em 59},~436--448.
\newblock {\url{https://doi.org/10.1111/aje.12846}}.

\bibitem[Mbumwae(1988)]{Mbumwae}
Mbumwae, L.L.
\newblock Environmental management of the Zambezi river system.
\newblock {\em Regul. Rivers Res. Manag.} {\bf 1988}, {\em
  2},~553--557.
\newblock {\url{https://doi.org/10.1002/rrr.3450020408}}.

\bibitem[Nehama and Reason(2021)]{NEHAMA2021101891}
Nehama, F.P.; Reason, C.J.
\newblock The wind-driven response of the Zambezi River plume along the Sofala
  Bank: A numerical model study.
\newblock {\em Reg. Stud. Mar. Sci.} {\bf 2021}, {\em
  46},~101891.
\newblock {\url{https://doi.org/10.1016/j.rsma.2021.101891}}.

\bibitem[Gope et~al.(2015)Gope, Sass-Klaassen, Irvine, Beevers, and Hes]{Gope}
Gope, E.T.; Sass-Klaassen, U.G.W.; Irvine, K.; Beevers, L.; Hes, E.M.A.
\newblock Effects of flow alteration on Apple-ring Acacia (Faidherbia albida)
  stands, Middle Zambezi floodplains, Zimbabwe.
\newblock {\em Ecohydrology} {\bf 2015}, {\em 8},~922--934.
\newblock {\url{https://doi.org/10.1002/eco.1541}}.

\bibitem[Moore et~al.(2022)Moore, Cotterill, Main, and Williams]{Moore}
Moore, A.E.; Cotterill, F.P.; Main, M.P.; Williams, H.B., The Zambezi: Origins
  and Legacies of Earth's Oldest River System.
\newblock In {\em Large Rivers}; John Wiley \& Sons, Ltd.: Hoboken, NJ, USA, 2022; Chapter~16, % Author's reply: yes, confirm.
  pp. 457--487.
\newblock {\url{https://doi.org/10.1002/9781119412632.ch16}}.

\bibitem[Leal et~al.(2009)Leal, Sá, Nordez, Brotas, and Paula]{LEAL2009605}
Leal, M.; Sá, C.; Nordez, S.; Brotas, V.; Paula, J.
\newblock Distribution and vertical dynamics of planktonic communities at
  Sofala Bank, Mozambique.
\newblock {\em Estuar. Coast. Shelf Sci.} {\bf 2009}, {\em
  84},~605--616.
\newblock {\url{https://doi.org/10.1016/j.ecss.2009.07.028}}.

\bibitem[Malauene et~al.(2018)Malauene, Moloney, Lett, Roberts, Marsac, and
  Penven]{MALAUENE20181}
Malauene, B.S.; Moloney, C.L.; Lett, C.; Roberts, M.J.; Marsac, F.; Penven, P.
\newblock Impact of offshore eddies on shelf circulation and river plumes of
  the Sofala Bank, Mozambique Channel.
\newblock {\em J. Mar. Syst.} {\bf 2018}, {\em 185},~1--12.
\newblock {\url{https://doi.org/10.1016/j.jmarsys.2018.05.001}}.

\bibitem[Malauene et~al.(2021)Malauene, Lett, Marsac, Roberts, Brito, Abdula,
  and Moloney]{MALAUENE2021107268}
Malauene, B.S.; Lett, C.; Marsac, F.; Roberts, M.J.; Brito, A.; Abdula, S.;
  Moloney, C.L.
\newblock Spawning areas of two shallow-water penaeid shrimps (Penaeus indicus
  and Metapenaeus monoceros) on the Sofala Bank, Mozambique.
\newblock {\em Estuar. Coast. Shelf Sci.} {\bf 2021}, {\em
  253},~107268.
\newblock {\url{https://doi.org/10.1016/j.ecss.2021.107268}}.

\bibitem[{Jokar Arsanjani} et~al.(2018){Jokar Arsanjani}, Fibæk, and
  Vaz]{JOKARARSANJANI201838}
{Jokar Arsanjani}, J.; Fibæk, C.S.; Vaz, E.
\newblock Development of a cellular automata model using open source
  technologies for monitoring urbanisation in the global south: The case of
  Maputo, Mozambique.
\newblock {\em Habitat Int.} {\bf 2018}, {\em 71},~38--48.
\newblock {\url{https://doi.org/10.1016/j.habitatint.2017.11.003}}.

\bibitem[de~Sousa et~al.(2006)de~Sousa, Brito, Abdula, and
  Caputi]{SOUSA2006207}
de~Sousa, L.P.; Brito, A.; Abdula, S.; Caputi, N.
\newblock Research assessment for the management of the industrial
  shallow-water multi-species shrimp fishery in Sofala Bank in Mozambique.
\newblock {\em Fish. Res.} {\bf 2006}, {\em 77},~207--219.
\newblock {\url{https://doi.org/10.1016/j.fishres.2005.10.009}}.

\bibitem[Miguel(2003)]{MIGUEL2003365}
Miguel, J.J.
\newblock On optimal choice of delay equations to model shrimp population
  dynamics in Sofala Bank, Mozambique.
\newblock {\em Nonlinear Anal. Real World Appl.} {\bf 2003}, {\em
  4},~365--371.
\newblock {\url{https://doi.org/10.1016/S1468-1218(02)00056-1}}.

\bibitem[Gammelsrød(1992)]{GAMMELSROD199291}
Gammelsrød, T.
\newblock Variation in shrimp abundance on the Sofala Bank, Mozambique, and its
  relation to the Zambezi River runoff.
\newblock {\em Estuar. Coast. Shelf Sci.} {\bf 1992}, {\em
  35},~91--103.
\newblock {\url{https://doi.org/10.1016/S0272-7714(05)80058-7}}.

\bibitem[Nhangumbe et~al.(2023)Nhangumbe, Nascetti, Georganos, and
  Ban]{NHANGUMBE2023101015}
Nhangumbe, M.; Nascetti, A.; Georganos, S.; Ban, Y.
\newblock Supervised and unsupervised machine learning approaches using
  Sentinel data for flood mapping and damage assessment in Mozambique.
\newblock {\em Remote Sens. Appl. Soc. Environ.} {\bf
  2023}, {\em 32},~101015.
\newblock {\url{https://doi.org/10.1016/j.rsase.2023.101015}}.

\bibitem[Salvucci and Santos(2020)]{SALVUCCI2020106713}
Salvucci, V.; Santos, R.
\newblock Vulnerability to Natural Shocks: Assessing the Short-Term Impact on
  Consumption and Poverty of the 2015 Flood in Mozambique.
\newblock {\em Ecol. Econ.} {\bf 2020}, {\em 176},~106713.
\newblock {\url{https://doi.org/10.1016/j.ecolecon.2020.106713}}.

\bibitem[Gall(2004)]{Gall}
Gall, M.
\newblock Where to Go? Strategic Modelling of Access to Emergency Shelters in
  Mozambique.
\newblock {\em Disasters} {\bf 2004}, {\em 28},~82--97.
\newblock {\url{https://doi.org/10.1111/j.0361-3666.2004.00244.x}}.

\bibitem[Come et~al.(2023)Come, Peer, Nhamussua, Miranda, Macamo, Cabral,
  Madivadua, Zacarias, Narciso, and Snow]{COME2023106813}
Come, J.; Peer, N.; Nhamussua, J.L.; Miranda, N.A.; Macamo, C.C.; Cabral, A.S.;
  Madivadua, H.; Zacarias, D.; Narciso, J.; Snow, B.
\newblock A socio-ecological survey in Inhambane Bay mangrove ecosystems:
  Biodiversity, livelihoods, and conservation.
\newblock {\em Ocean Coast. Manag.} {\bf 2023}, {\em 244},~106813.
\newblock {\url{https://doi.org/10.1016/j.ocecoaman.2023.106813}}.

\bibitem[Fichtner et~al.(2023)Fichtner, Mandery, Wieland, Groth, Martinis, and
  Riedlinger]{FICHTNER2023103329}
Fichtner, F.; Mandery, N.; Wieland, M.; Groth, S.; Martinis, S.; Riedlinger, T.
\newblock Time-series analysis of Sentinel-1/2 data for flood detection using a
  discrete global grid system and seasonal decomposition.
\newblock {\em Int. J. Appl. Earth Obs. Geoinf.} {\bf 2023}, {\em 119},~103329.
\newblock {\url{https://doi.org/10.1016/j.jag.2023.103329}}.

\bibitem[Bofana et~al.(2022)Bofana, Zhang, Wu, Zeng, Nabil, Zhang, Elnashar,
  Tian, {da Silva}, Botão, Atumane, Mushore, and Yan]{BOFANA2022112808}
Bofana, J.; Zhang, M.; Wu, B.; Zeng, H.; Nabil, M.; Zhang, N.; Elnashar, A.;
  Tian, F.; {da Silva}, J.M.; Botão, A.;  et~al.
\newblock How long did crops survive from floods caused by Cyclone Idai in
  Mozambique detected with multi-satellite data.
\newblock {\em Remote Sens. Environ.} {\bf 2022}, {\em 269},~112808.
\newblock {\url{https://doi.org/10.1016/j.rse.2021.112808}}.

\bibitem[Guldemond and van Aarde(2010)]{Guldemond}
Guldemond, R.A.R.; van Aarde, R.J.
\newblock Forest patch size and isolation as drivers of bird species richness
  in Maputaland, Mozambique.
\newblock {\em J. Biogeogr.} {\bf 2010}, {\em 37},~1884--1893.
\newblock {\url{https://doi.org/10.1111/j.1365-2699.2010.02338.x}}.

\bibitem[Ramayanti et~al.(2022)Ramayanti, Nur, Syifa, Panahi, Achmad, Park, and
  Lee]{RAMAYANTI20221025}
Ramayanti, S.; Nur, A.S.; Syifa, M.; Panahi, M.; Achmad, A.R.; Park, S.; Lee,
  C.W.
\newblock Performance comparison of two deep learning models for flood
  susceptibility map in Beira area, Mozambique.
\newblock {\em  Egypt. J. Remote Sens. Space Sci.} {\bf
  2022}, {\em 25},~1025--1036.
\newblock {\url{https://doi.org/10.1016/j.ejrs.2022.11.003}}.

\bibitem[Smith et~al.(2019)Smith, Ryan, Vollmer, Woollen, Keane, Fisher,
  Baumert, Grundy, Carvalho, Lisboa, Luz, Zorrilla-Miras, Patenaude, Ribeiro,
  Artur, and Mahamane]{SMITH2019101976}
Smith, H.E.; Ryan, C.M.; Vollmer, F.; Woollen, E.; Keane, A.; Fisher, J.A.;
  Baumert, S.; Grundy, I.M.; Carvalho, M.; Lisboa, S.N.;  et~al.
\newblock Impacts of land use intensification on human wellbeing: Evidence from
  rural Mozambique.
\newblock {\em Glob. Environ. Chang.} {\bf 2019}, {\em 59},~101976.
\newblock {\url{https://doi.org/10.1016/j.gloenvcha.2019.101976}}.

\bibitem[Lundgren and Strandh(2022)]{LUNDGREN2022103023}
Lundgren, M.; Strandh, V.
\newblock Navigating a double burden – Floods and social vulnerability in
  local communities in rural Mozambique.
\newblock {\em Int. J. Disaster Risk Reduct.} {\bf 2022},
  {\em 77},~103023.
\newblock {\url{https://doi.org/10.1016/j.ijdrr.2022.103023}}.

\bibitem[Silva et~al.(2018)Silva, Loboda, and Strong]{SILVA201843}
Silva, J.A.; Loboda, T.; Strong, M.
\newblock Examining aspiration’s imprint on the landscape: Lessons from
  Mozambique’s Limpopo National Park.
\newblock {\em Glob. Environ. Chang.} {\bf 2018}, {\em 51},~43--53.
\newblock {\url{https://doi.org/10.1016/j.gloenvcha.2018.04.013}}.

\bibitem[Martins and Shackleton(2022)]{MARTINS2022102793}
Martins, A.R.; Shackleton, C.M.
\newblock The contribution of wild palms to the livelihoods and diversification
  of rural households in southern Mozambique.
\newblock {\em For. Policy Econ.} {\bf 2022}, {\em 142},~102793.
\newblock {\url{https://doi.org/10.1016/j.forpol.2022.102793}}.

\bibitem[Pittman et~al.(2023)Pittman, Swanborn, Connor, and
  Wright]{PITTMAN2023}
Pittman, S.J.; Swanborn, D.J.; Connor, D.W.; Wright, D.J.
\newblock Application of Estuarine and Coastal Classifications in Marine
  Spatial Management. In {\em Reference Module in Earth Systems and
  Environmental Sciences}; Elsevier: Amsterdam, The Netherlands, 2023. % Author's reply: yes, confirm.
\newblock {\url{https://doi.org/10.1016/B978-0-323-90798-9.00040-8}}.

\bibitem[Xi et~al.(2019)Xi, Du, Wang, and Zhang]{XI2019111212}
Xi, W.; Du, S.; Wang, Y.C.; Zhang, X.
\newblock A spatiotemporal cube model for analyzing satellite image time
  series: Application to land-cover mapping and change detection.
\newblock {\em Remote Sens. Environ.} {\bf 2019}, {\em 231},~111212.
\newblock {\url{https://doi.org/10.1016/j.rse.2019.111212}}.

\bibitem[Abidi et~al.(2023)Abidi, Ienco, Abbes, and Farah]{ABIDI2023106152}
Abidi, A.; Ienco, D.; Abbes, A.B.; Farah, I.R.
\newblock Combining 2D encoding and convolutional neural network to enhance
  land cover mapping from Satellite Image Time Series.
\newblock {\em Eng. Appl. Artif. Intell.} {\bf 2023},
  {\em 122},~106152.
\newblock {\url{https://doi.org/10.1016/j.engappai.2023.106152}}.

\bibitem[Mohammadi et~al.(2023)Mohammadi, Belgiu, and Stein]{MOHAMMADI2023272}
Mohammadi, S.; Belgiu, M.; Stein, A.
\newblock Improvement in crop mapping from satellite image time series by
  effectively supervising deep neural networks.
\newblock {\em ISPRS J. Photogramm. Remote Sens.} {\bf 2023},
  {\em 198},~272--283.
\newblock {\url{https://doi.org/10.1016/j.isprsjprs.2023.03.007}}.

\bibitem[Namaganda et~al.(2023)Namaganda, Otsuki, and
  Steel]{NAMAGANDA2023117811}
Namaganda, E.; Otsuki, K.; Steel, G.
\newblock Understanding the cumulative socioenvironmental impacts of energy
  transition-induced extractivism in Mozambique: The role of mixed methods.
\newblock {\em J. Environ. Manag.} {\bf 2023}, {\em
  338},~117811.
\newblock {\url{https://doi.org/10.1016/j.jenvman.2023.117811}}.

\bibitem[Gemusse et~al.(2023)Gemusse, Cardoso-Fernandes, Lima, and
  Teodoro]{GEMUSSE2023101022}
Gemusse, U.; Cardoso-Fernandes, J.; Lima, A.; Teodoro, A.
\newblock Identification of pegmatites zones in Muiane and Naipa (Mozambique)
  from Sentinel-2 images, using band combinations, band ratios, PCA and
  supervised classification.
\newblock {\em Remote Sens. Appl. Soc. Environ.} {\bf
  2023}, {\em 32},~101022.
\newblock {\url{https://doi.org/10.1016/j.rsase.2023.101022}}.

\bibitem[Liang et~al.(2023)Liang, Duncanson, Silva, and
  Sedano]{LIANG2023113367}
Liang, M.; Duncanson, L.; Silva, J.A.; Sedano, F.
\newblock Quantifying aboveground biomass dynamics from charcoal degradation in
  Mozambique using GEDI Lidar and Landsat.
\newblock {\em Remote Sens. Environ.} {\bf 2023}, {\em 284},~113367.
\newblock {\url{https://doi.org/10.1016/j.rse.2022.113367}}.

\bibitem[Bey et~al.(2020)Bey, Jetimane, Lisboa, Ribeiro, Sitoe, and
  Meyfroidt]{BEY2020111611}
Bey, A.; Jetimane, J.; Lisboa, S.N.; Ribeiro, N.; Sitoe, A.; Meyfroidt, P.
\newblock Mapping smallholder and large-scale cropland dynamics with a flexible
  classification system and pixel-based composites in an emerging frontier of
  Mozambique.
\newblock {\em Remote Sens. Environ.} {\bf 2020}, {\em 239},~111611.
\newblock {\url{https://doi.org/10.1016/j.rse.2019.111611}}.

\bibitem[Lemenkova(2022)]{Lemenkova202213}
Lemenkova, P.
\newblock {Mapping submarine geomorphology of the Philippine and Mariana
  trenches by an automated approach using GMT scripts}.
\newblock {\em Proc. Latv. Acad. Sciences. Sect. B Nat. Exact Appl. Sci.} {\bf 2022}, {\em 76},~258--266.
\newblock {\url{https://doi.org/10.2478/prolas-2022-0039}}.

\bibitem[Lemenkova and Debeir(2023)]{jimaging9050098}
Lemenkova, P.; Debeir, O.
\newblock Multispectral Satellite Image Analysis for Computing Vegetation
  Indices by R in the Khartoum Region of Sudan, Northeast Africa.
\newblock {\em J. Imaging} {\bf 2023}, {\em 9}, 98.
\newblock {\url{https://doi.org/10.3390/jimaging9050098}}.

\bibitem[Lemenkova and Debeir(2022)]{Lemenkova202229}
Lemenkova, P.; Debeir, O.
\newblock {R Libraries for Remote Sensing Data Classification by k-means
  Clustering and NDVI Computation in Congo River Basin, DRC}.
\newblock {\em Appl. Sci.} {\bf 2022}, {\em 12},~12554.
\newblock {\url{https://doi.org/10.3390/app122412554}}.

\bibitem[Lemenkova and Debeir(2023{\natexlab{a}})]{technologies11020046}
Lemenkova, P.; Debeir, O.
\newblock GDAL and PROJ Libraries Integrated with GRASS GIS for Terrain
  Modelling of the Georeferenced Raster Image.
\newblock {\em Technologies} {\bf 2023}, {\em 11}, 46.
\newblock {\url{https://doi.org/10.3390/technologies11020046}}.

\bibitem[Lemenkova and Debeir(2023{\natexlab{b}})]{jmse11040871}
Lemenkova, P.; Debeir, O.
\newblock Computing Vegetation Indices from the Satellite Images Using GRASS
  GIS Scripts for Monitoring Mangrove Forests in the Coastal Landscapes of
  Niger Delta, Nigeria.
\newblock {\em J. Mar. Sci. Eng.} {\bf 2023}, {\em 11}, 871.
\newblock {\url{https://doi.org/10.3390/jmse11040871}}.

\bibitem[Wessel et~al.(2019)Wessel, Luis, Uieda, Scharroo, Wobbe, Smith, and
  Tian]{Wessel}
Wessel, P.; Luis, J.F.; Uieda, L.; Scharroo, R.; Wobbe, F.; Smith, W.H.F.;
  Tian, D.
\newblock The Generic Mapping Tools Version 6.
\newblock {\em Geochem. Geophys. Geosystems} {\bf 2019}, {\em
  20},~5556--5564.
\newblock {\url{https://doi.org/10.1029/2019GC008515}}.

\bibitem[Lemenkova(2022{\natexlab{a}})]{Lemenkova202217}
Lemenkova, P.
\newblock {Mapping Climate Parameters over the Territory of Botswana Using GMT
  and Gridded Surface Data from TerraClimate}.
\newblock {\em ISPRS Int. J. Geo-Inf.} {\bf 2022}, {\em
  11},~473.
\newblock {\url{https://doi.org/10.3390/ijgi11090473}}.

\bibitem[Lemenkova(2022{\natexlab{b}})]{Lemenkova202212}
Lemenkova, P.
\newblock {Handling Dataset with Geophysical and Geological Variables on the
  Bolivian Andes by the GMT Scripts}.
\newblock {\em Data} {\bf 2022}, {\em 7},~74.
\newblock {\url{https://doi.org/10.3390/data7060074}}.

\bibitem[Lemenkova(2022{\natexlab{c}})]{Lemenkova202203}
Lemenkova, P.
\newblock {Console-Based Mapping of Mongolia Using GMT Cartographic Scripting
  Toolset for Processing TerraClimate Data}.
\newblock {\em Geosciences} {\bf 2022}, {\em 12},~140.
\newblock {\url{https://doi.org/10.3390/geosciences12030140}}.

\bibitem[Lemenkova(2022{\natexlab{d}})]{Lemenkova202215}
Lemenkova, P.
\newblock {Cartographic scripts for seismic and geophysical mapping of
  Ecuador}.
\newblock {\em Geografie} {\bf 2022}, {\em 127},~195--218.
\newblock {\url{https://doi.org/10.37040/geografie.2022.006}}.

\bibitem[Mitchard et~al.(2009)Mitchard, Saatchi, Woodhouse, Nangendo, Ribeiro,
  Williams, Ryan, Lewis, Feldpausch, and Meir]{Mitchard}
Mitchard, E.T.A.; Saatchi, S.S.; Woodhouse, I.H.; Nangendo, G.; Ribeiro, N.S.;
  Williams, M.; Ryan, C.M.; Lewis, S.L.; Feldpausch, T.R.; Meir, P.
\newblock Using satellite radar backscatter to predict above-ground woody
  biomass: A consistent relationship across four different African landscapes.
\newblock {\em Geophys. Res. Lett.} {\bf 2009}, {\em 36%MDPI: Please add the page number.
},~1--6.
\newblock {\url{https://doi.org/10.1029/2009GL040692}}.

\bibitem[Fernando and Senanayake(2023)]{FERNANDO2023}
Fernando, W.A.M.; Senanayake, I.
\newblock Developing a two-decadal time-record of rice field maps using
  Landsat-derived multi-index image collections with a random forest
  classifier: A Google Earth Engine based approach.
\newblock {\em Inf. Process. Agric.} {\bf 2023}, \emph{in press}.
\newblock {\url{https://doi.org/10.1016/j.inpa.2023.02.009}}.

\bibitem[Kolarik et~al.(2024)Kolarik, Shrestha, Caughlin, and
  Brandt]{KOLARIK2024111445}
Kolarik, N.; Shrestha, N.; Caughlin, T.; Brandt, J.
\newblock Leveraging high resolution classifications and random forests for
  hindcasting decades of mesic ecosystem dynamics in the Landsat time series.
\newblock {\em Ecol. Indic.} {\bf 2024}, {\em 158},~111445.
\newblock {\url{https://doi.org/10.1016/j.ecolind.2023.111445}}.

\bibitem[Bera et~al.(2023)Bera, {Das Chatterjee}, Bera, Ghosh, and
  Dinda]{BERA20234691}
Bera, D.; {Das Chatterjee}, N.; Bera, S.; Ghosh, S.; Dinda, S.
\newblock Comparative performance of Sentinel-2 MSI and Landsat-8 OLI data in
  canopy cover prediction using Random Forest model: Comparing model
  performance and tuning parameters.
\newblock {\em Adv. Space Res.} {\bf 2023}, {\em 71},~4691--4709.
\newblock {\url{https://doi.org/10.1016/j.asr.2023.01.027}}.

\bibitem[Yan et~al.(2023)Yan, Li, Smith, Yang, Ma, and Su]{land12122149}
Yan, X.; Li, J.; Smith, A.R.; Yang, D.; Ma, T.; Su, Y.
\newblock Rapid Land Cover Classification Using a 36-Year Time Series of
  Multi-Source Remote Sensing Data.
\newblock {\em Land} {\bf 2023}, {\em 12}, 2149.
\newblock {\url{https://doi.org/10.3390/land12122149}}.

\bibitem[Zhang et~al.(2023)Zhang, Zhang, Liu, Lan, Gao, and
  Li]{ZHANG2023108250}
Zhang, H.; Zhang, Y.; Liu, K.; Lan, S.; Gao, T.; Li, M.
\newblock Winter wheat yield prediction using integrated Landsat 8 and
  Sentinel-2 vegetation index time-series data and machine learning algorithms.
\newblock {\em Comput. Electron. Agric.} {\bf 2023}, {\em
  213},~108250.
\newblock {\url{https://doi.org/10.1016/j.compag.2023.108250}}.

\bibitem[Mountrakis and Heydari(2023)]{MOUNTRAKIS2023106}
Mountrakis, G.; Heydari, S.S.
\newblock Harvesting the Landsat archive for land cover land use classification
  using deep neural networks: Comparison with traditional classifiers and
  multi-sensor benefits.
\newblock {\em ISPRS J. Photogramm. Remote Sens.} {\bf 2023},
  {\em 200},~106--119.
\newblock {\url{https://doi.org/10.1016/j.isprsjprs.2023.05.005}}.

\bibitem[Singh and Tyagi(2021)]{SINGH2021100645}
Singh, M.; Tyagi, K.D.
\newblock Pixel based classification for Landsat 8 OLI multispectral satellite
  images using deep learning neural network.
\newblock {\em Remote Sens. Appl. Soc. Environ.} {\bf
  2021}, {\em 24},~100645.
\newblock {\url{https://doi.org/10.1016/j.rsase.2021.100645}}.

\bibitem[Zheng et~al.(2023)Zheng, Yang, Liu, Xia, and Zhang]{ZHENG2023100974}
Zheng, Z.; Yang, B.; Liu, S.; Xia, J.; Zhang, X.
\newblock Extraction of impervious surface with Landsat based on machine
  learning in Chengdu urban, China.
\newblock {\em Remote Sens. Appl. Soc. Environ.} {\bf
  2023}, {\em 30},~100974.
\newblock {\url{https://doi.org/10.1016/j.rsase.2023.100974}}.

\bibitem[{Brown de Colstoun} et~al.(2003){Brown de Colstoun}, Story, Thompson,
  Commisso, Smith, and Irons]{BROWNDECOLSTOUN2003316}
{Brown de Colstoun}, E.C.; Story, M.H.; Thompson, C.; Commisso, K.; Smith,
  T.G.; Irons, J.R.
\newblock National Park vegetation mapping using multitemporal Landsat 7 data
  and a decision tree classifier.
\newblock {\em Remote Sens. Environ.} {\bf 2003}, {\em 85},~316--327.
\newblock {\url{https://doi.org/10.1016/S0034-4257(03)00010-5}}.

\bibitem[Vaidya et~al.(2024)Vaidya, Keskar, and Kotharkar]{VAIDYA2024101770}
Vaidya, M.; Keskar, R.; Kotharkar, R.
\newblock Classifying heterogeneous urban form into local climate zones using
  supervised learning and greedy clustering incorporating Landsat dataset.
\newblock {\em Urban Clim.} {\bf 2024}, {\em 53},~101770.
\newblock {\url{https://doi.org/10.1016/j.uclim.2023.101770}}.

\bibitem[Pasquarella et~al.(2023)Pasquarella, Morreale, Brown, Kilbride, and
  Thompson]{PASQUARELLA2023103561}
Pasquarella, V.J.; Morreale, L.L.; Brown, C.F.; Kilbride, J.B.; Thompson, J.R.
\newblock Not-so-random forests: Comparing voting and decision tree ensembles
  for characterizing partial harvest events.
\newblock {\em Int. J. Appl. Earth Obs. Geoinf.} {\bf 2023}, {\em 125},~103561.
\newblock {\url{https://doi.org/10.1016/j.jag.2023.103561}}.

\bibitem[Al-Obeidat et~al.(2015)Al-Obeidat, Al-Taani, Belacel, Feltrin, and
  Banerjee]{ALOBEIDAT20151192}
Al-Obeidat, F.; Al-Taani, A.T.; Belacel, N.; Feltrin, L.; Banerjee, N.
\newblock A Fuzzy Decision Tree for Processing Satellite Images and Landsat
  Data.
\newblock {\em Procedia Comput. Sci.} {\bf 2015}, {\em 52},~1192--1197.
\newblock {\url{https://doi.org/10.1016/j.procs.2015.05.157}}.

\end{thebibliography}

%\PublishersNote{}

%%%%%%%%%%%%%%%%%%%%%%%%%%%%%%%%%%%%%%%%%%
\end{adjustwidth}
\end{document}
